\documentclass[12pt, titlepage]{article}

\usepackage{amsmath, mathtools}

\usepackage[round]{natbib}
\usepackage{amsfonts}
\usepackage{amssymb}
\usepackage{graphicx}
\usepackage{colortbl}
\usepackage{xr}
\usepackage{hyperref}
\usepackage{longtable}
\usepackage{xfrac}
\usepackage{tabularx}
\usepackage{float}
\usepackage{siunitx}
\usepackage{booktabs}
\usepackage{multirow}
\usepackage[section]{placeins}
\usepackage{caption}
\usepackage{fullpage}
\usepackage{parskip}

\hypersetup{
bookmarks=true,     % show bookmarks bar?
colorlinks=true,       % false: boxed links; true: colored links
linkcolor=red,          % color of internal links (change box color with linkbordercolor)
citecolor=blue,      % color of links to bibliography
filecolor=magenta,  % color of file links
urlcolor=cyan          % color of external links
}

\usepackage{array}

\externaldocument{../../SRS/SRS}

\newcommand{\descriptiveid}[1]{\texttt{\detokenize{#1}}}
% Link into the published SRS PDF. Works in a browser/PDF viewer and on
% GitHub, unlike a relative-path link. Usage: \srsref{FR_AssignTA}
\newcommand{\srsurl}{https://ayushpatel2026.github.io/Keystone/SRS/SRS.pdf}
\newcommand{\srsref}[1]{\href{\srsurl\#nameddest=#1}{\descriptiveid{#1}}}

\input{../../Comments.text}
\input{../../Common.text}

\newcommand{\moduleid}[1]{\texttt{\detokenize{#1}}}

\begin{document}

\title{Module Interface Specification for \projname{}}

\author{\authname}

\date{\today}

\maketitle

\pagenumbering{roman}

\section{Revision History}

\begin{tabularx}{\textwidth}{p{3cm}p{2cm}X}
\toprule {\bf Date} & {\bf Version} & {\bf Notes}\\
\midrule
Nov. 15, 2026 & 1.0 & Added the MIS of the Projects, Students and Teams, Recommendation, Approval Workflow, and Threads modules, and the shared types they need\\
\bottomrule
\end{tabularx}

~\newpage

\section{Symbols, Abbreviations and Acronyms}

See the SRS documentation at
\url{https://github.com/Ayushpatel2026/Keystone/raw/main/pdfs/SRS/SRS.pdf}.
The symbols used to write the specifications are listed in
Table~\ref{TblNotation}.

\newpage

\tableofcontents

\newpage

\pagenumbering{arabic}

\section{Introduction}

The following document details the Module Interface Specifications for
\projname{}, a platform that makes capstone project proposals, project discovery,
team formation, and team project approval transparent, consistent, and easy to
manage for students, supervisors, and course instructors.

Complementary documents include the System Requirement Specifications
and Module Guide.  The full documentation and implementation can be
found at \url{https://github.com/Ayushpatel2026/Keystone}.

\section{Notation}
\label{SecNotation}

The module specifications follow \citet{HoffmanAndStrooper1995}, with template
module conventions adapted from \citet{GhezziEtAl2003}. Each access program
specifies its inputs, output, state transition, and exception conditions.

\begin{table}[H]
\centering
\renewcommand{\arraystretch}{1.15}
\begin{tabularx}{\textwidth}{l X}
\toprule
\textbf{Notation} & \textbf{Meaning} \\
\midrule
bool & Boolean values $\{\text{true},\text{false}\}$. \\
String & Finite sequence of characters where $\epsilon$ is the empty string. \\
$\mathbb{N}$ & Natural numbers $\{0,1,2,\dots\}$. \\
$[0,1]$ & Real numbers from 0 to 1 inclusive. \\
$\emptyset$ & Empty set or map, according to its declared type. \\
None & Absence of a value where $T\cup\{\text{None}\}$ denotes an optional value. \\
$x$, $x'$ & Values of state variable $x$ before and after a routine. \\
$\in$, $\notin$ & Membership and non-membership in a set. \\
$\ne$ & Not equal. \\
$\land$, $\neg$ & Logical and, and logical negation (not). \\
$P\Rightarrow Q$ & If condition $P$ holds, then property or transition $Q$ holds. \\
$out$, $exc$ & Routine return value and raised exception, respectively. \\
$\mathcal{V}$, $\mathcal{F}$ & Sets of registered UI view and form names. \\
$\operatorname{Map}\langle K,V\rangle$ & Mapping from keys of type $K$ to values of type $V$. \\
$\operatorname{Set}\langle T\rangle$ & Set of values of type $T$. \\
$\operatorname{List}\langle T\rangle$ & Ordered sequence of values of type $T$. \\
$\vert S\vert$ & Number of elements in set or list $S$. \\
$k\mapsto v$ & Associates key $k$ with value $v$. \\
$\operatorname{dom}(M)$ & Set of keys in mapping $M$. \\
$M[k]$ & Value associated with key $k$ in mapping $M$. \\
$M\cup\{k\mapsto v\}$ & Adds a new key-value pair to mapping $M$. \\
$M\setminus\{k\mapsto v\}$ & Removes a key-value pair from mapping $M$. \\
$M\oplus\{k\mapsto v\}$ & Adds or replaces the value associated with key $k$ in mapping $M$. \\
-- & No input, program return value, or exported exception, according to column. \\
\bottomrule
\end{tabularx}
\caption{Notation used in the module specifications}
\label{TblNotation}
\end{table}

A transition describes the post-state of a successful call. State variables not mentioned remain unchanged. Unless stated otherwise, an exported exception leaves module state unchanged and produces no normal output. Output denotes a program return value, effects on the screen and other external resources are stated separately.

\newpage

\section{Module Decomposition}

Table~\ref{TblMH} follows the module hierarchy in the Module Guide~\citep{MG},
with each module's full name and identifier shown.

\begin{table}[h!]
\centering
\begin{tabular}{p{0.3\textwidth} p{0.6\textwidth}}
\toprule
\textbf{Level 1} & \textbf{Level 2}\\
\midrule

{Hardware-Hiding Module} & Hardware-Hiding Module (\moduleid{M_HardwareHiding}) \\
\midrule

\multirow{13}{0.3\textwidth}{Behaviour-Hiding Module} & User Interface Module (\moduleid{M_UserInterface})\\
 & Authentication Module (\moduleid{M_Authentication})\\
 & Authorization Module (\moduleid{M_Authorization})\\
 & Course Management Module (\moduleid{M_CourseManagement})\\
 & Enrollment Module (\moduleid{M_Enrollment})\\
 & Projects Module (\moduleid{M_Projects})\\
 & Students and Teams Module (\moduleid{M_StudentsTeams})\\
 & Recommendation Module (\moduleid{M_Recommendation})\\
 & Approval Workflow Module (\moduleid{M_ApprovalWorkflow})\\
 & Threads Module (\moduleid{M_Threads})\\
 & Notifications Module (\moduleid{M_Notifications})\\
 & Checklist Module (\moduleid{M_Checklist})\\
 & Course Report Module (\moduleid{M_CourseReport})\\
\midrule

\multirow{2}{0.3\textwidth}{Software Decision Module} & Persistence Module (\moduleid{M_Persistence})\\
 & Audit Log Module (\moduleid{M_AuditLog})\\
\bottomrule

\end{tabular}
\caption{Module Hierarchy}
\label{TblMH}
\end{table}

\newpage


\section{Data Types}
\label{SecDataTypes}

The following data types are shared across the Keystone modules.
\subsection{Enums}

\subsubsection{Action}
\begin{enumerate}
    \item Create, Read, Update, Delete: Create, view, modify, or remove a record.
    \item Submit, Apply, Withdraw: Submit a record or application, apply to a project, or withdraw it.
    \item Approve, Reject, RequestChanges: Record an approval decision or request revisions.
    \item Accept, Decline: Accept or decline a project or agreement.
    \item Import, ConfirmImport, Export: Import, confirm an enrollment import, or export course information.
    \item Assign, Verify: Assign a course role or verify an external supervisor.
    \item Configure, Enable, Disable: Change course settings or enable or disable a feature.
    \item AddMember, RemoveMember, LeaveTeam, DissolveTeam: Manage team membership.
    \item ChangeCredentials, ChangeEmail, LinkAuthentication, ChangeCourseAssociation: Change account credentials, email, linked sign-in, or course association.
    \item SendMessage, ViewReport: Send a discussion message or view a course report.
\end{enumerate}

\subsubsection{AuthProvider}
\begin{enumerate}
    \item McMasterSSO: Authentication provider for McMaster users.
    \item ExternalProvider: Trusted authentication provider for external supervisors.
\end{enumerate}

\subsubsection{Feature}
\begin{enumerate}
    \item TeamFormation: Team formation functionality.
    \item ProjectApprovals: Project approval functionality.
    \item ProjectProposals: Project proposal functionality.
\end{enumerate}

\subsubsection{SupervisorVerification}
\begin{enumerate}
    \item Unverified: External supervisor account state at sign-in until an instructor approves that user's project proposal.
    \item Verified: External supervisor account state after an instructor approves that user's project proposal.
\end{enumerate}


\subsubsection{EnrollmentStatus}
\begin{enumerate}
    \item Enrolled: Student is enrolled in the course.
    \item Flagged: Student is absent from the latest imported enrollment list and requires review.
\end{enumerate}

\subsubsection{EnrollmentSource}
\begin{enumerate}
    \item Avenue: Enrollment data obtained from Avenue to Learn.
    \item UploadedFile: Enrollment data provided through a file.
\end{enumerate}

\subsubsection{UserRole}
\begin{enumerate}
    \item Student: A student participating in a capstone course.
    \item Supervisor: A user proposing or supervising projects.
    \item Instructor: A user responsible for managing a capstone course.
    \item TA: A teaching assistant supporting course management.
\end{enumerate}

\subsubsection{ProjectStatus}
\begin{enumerate}
    \item Draft: Project proposal has not been submitted.
    \item ProposalSubmitted: Proposal is awaiting review.
    \item ProposalChangesRequested: Modifications are required.
    \item ProposalRejected: Project proposal was rejected.
    \item ProposalWithdrawn: Project proposal was withdrawn.
    \item Open: Listed project has available capacity.
    \item Full: Listed project has no remaining capacity.
    \item Closed: Project selection period has ended.
\end{enumerate}

\subsubsection{TeamProjectStatus}
\begin{enumerate}
    \item Applied: Team has applied to a project.
    \item SupervisorAccepts: Supervisor has accepted the team.
    \item Agreed: Team has accepted the working arrangement.
    \item TeamProjectSubmitted: Team project was submitted for approval.
    \item TeamProjectUnderReview: Submission is under review.
    \item TeamProjectChangesRequested: Modifications are required.
    \item TeamProjectRejected: Team project was rejected.
    \item Approved: Team project has been approved.
    \item ApprovedChangeUnderReview: Changes to an approved
    project are under review.
    \item DeclinedOrWithdrawn: Application was declined or withdrawn.
\end{enumerate}

\subsubsection{ApprovalMode}
\begin{enumerate}
    \item InPlatform: Approval is performed within Keystone.
    \item External: Approval occurs externally and is recorded manually.
    \item Integrated: Approval is synchronized through an external adapter.
\end{enumerate}

\subsubsection{Program}
\begin{enumerate}
    \item SE, CS, TRON: The program of a student, as in the SRS.
\end{enumerate}

\subsubsection{TeamMemberRole}
\begin{enumerate}
    \item PointOfContact: Team member who is the team's contact for supervisors and staff, and who manages membership.
    \item Member: Ordinary team member.
\end{enumerate}

\subsubsection{OfferKind}
\begin{enumerate}
    \item Invitation: A team invites a student to join.
    \item Request: A student asks to join a team.
\end{enumerate}

\subsubsection{RuleViolation}
\begin{enumerate}
    \item TooFewMembers: Team has fewer members than the course minimum.
    \item TooManyMembers: Team has more members than the course maximum.
    \item ProgramCombinationNotAllowed: The programs of the members are not a permitted combination.
\end{enumerate}

\subsubsection{RequestKind}
\begin{enumerate}
    \item Initial: First in-platform approval request for a team's project.
    \item Revision: In-platform request to change a project that is already approved.
\end{enumerate}

\subsubsection{ThreadTarget}
\begin{enumerate}
    \item ProjectProposal: A discussion attached to a project or its proposal.
    \item Application: A discussion attached to a team's project application.
    \item Team: A discussion attached to a team's formation.
\end{enumerate}

\subsection{Data Structures}

\subsubsection{Dictionary}
For types $K$ and $V$, a finite mapping from keys of type $K$ to values of type $V$, written $\operatorname{Map}\langle K,V\rangle$.

\subsubsection{List}
For any type $T$, an ordered collection of elements of type $T$.

\subsubsection{DateTime}
A date and time, as in the SRS data dictionary.

\subsubsection{CourseResource}
A course document or link with an identifier, course association, and publication status.

\subsubsection{SessionID}
A string that uniquely identifies an authenticated session.

\subsubsection{AuthProof}
Credentials, a verification code, or an identity provider response submitted to verify a user's identity.

\subsubsection{FieldValues}
\[
\text{FieldValues} = \operatorname{Map}\langle\text{string},\text{Value}\rangle
\]
Maps form field names to entered values.

\subsubsection{FieldErrors}
\[
\text{FieldErrors} = \operatorname{Map}\langle\text{string},\text{string}\rangle
\]
Maps form field names to error messages.

\subsubsection{EnrollmentEntry}
A record containing a student's name, email or MacID, program, and associated course. These are the minimum fields required for enrollment import by the SRS.

\subsubsection{EnrollmentPreview}
A record containing the proposed enrollment additions, existing matches, students absent from the imported list, validation errors, and an identifier for the pending import.

\subsubsection{ImportID}
A unique string identifying a pending enrollment import.

\subsubsection{ViewData}
The record or collection supplied by the owning module to populate a view, with fields and types defined in that module's specification.

\subsubsection{User}
A record representing a Keystone user, including their identity and assigned roles.

\subsubsection{Course}
A record containing capstone course information, configuration settings, and enrollment information. The course settings used by the Students and Teams, Approval Workflow, and Projects modules are \texttt{minTeamSize} ($\mathbb{N}$), \texttt{maxTeamSize} ($\mathbb{N}$), \texttt{allowedCombinations} ($\operatorname{Set}\langle\operatorname{Set}\langle\text{Program}\rangle\rangle$), \texttt{approvalMode} (ApprovalMode), \texttt{externalAdapter} (the configuration of the external approval system), and \texttt{submissionRequirements} ($\operatorname{Map}\langle\text{String},\text{bool}\rangle$, each submission item mapped to true if it is required).

\subsubsection{StudentProfile}
A record containing a student's background, skills, interests, experience, and project preferences. It also contains the student's \emph{program} (Program) and \emph{seekingTeam} (bool, true if the student is looking for a team).

\subsubsection{Team}
A record representing a group of students participating in the capstone project selection process, with fields $\mathit{id}$, $\mathit{course}$, $\mathit{members}$ ($\operatorname{Map}\langle\text{User},\text{TeamMemberRole}\rangle$), $\mathit{membersNeeded}$ ($\mathbb{N}$, the number of members the team is recruiting), $\mathit{ta}$ ($\text{User}\cup\{\text{None}\}$), and $\mathit{profile}$ (TeamProfile).

\subsubsection{TeamProfile}
A record containing a team's collective skills, interests, goals, and work style.

\subsubsection{Project}
A record containing project information, supervisor capacity, course association, and ProjectStatus, with fields $\mathit{id}$, $\mathit{course}$, $\mathit{proposer}$ (User), $\mathit{supervisor}$ ($\text{User}\cup\{\text{None}\}$), $\mathit{team}$ ($\text{Team}\cup\{\text{None}\}$, the team that proposed a student-proposed project), $\mathit{fields}$ (FieldValues), $\mathit{status}$ (ProjectStatus), $\mathit{capacity}$ ($\mathbb{N}$), $\mathit{reserved}$ ($\operatorname{Set}\langle\text{String}\rangle$, the identifiers of the teams holding a slot), and $\mathit{comment}$ (String, the reviewer's feedback).

\subsubsection{ProjectApplication}
A record representing a team's application to a project, with fields $\mathit{id}$, $\mathit{project}$, $\mathit{team}$, and the associated TeamProjectStatus $\mathit{status}$.

\subsubsection{Offer}
A record of a pending invitation or join request, with fields $\mathit{id}$, $\mathit{team}$, $\mathit{user}$ (the student invited or asking), and $\mathit{kind}$ (OfferKind).

\subsubsection{ApprovalRequest}
A record of an in-platform request to approve a team project or a change to it, with fields $\mathit{id}$, $\mathit{application}$, $\mathit{content}$ (FieldValues), $\mathit{kind}$ (RequestKind), $\mathit{status}$ (TeamProjectStatus), and $\mathit{comment}$ (String).

\subsubsection{SyncStatus}
A record with fields $\mathit{lastSuccess}$ ($\text{DateTime}\cup\{\text{None}\}$) and $\mathit{failed}$ (bool), describing the latest attempt to read approval outcomes from a course's external system.

\subsubsection{Notification}
A record containing information about an event or action communicated to a user.

\subsubsection{ChecklistItem}
A record describing a required user action and its completion status.

\subsubsection{Thread}
A record representing a discussion attached to a project proposal, a project application, or a team formation, with fields $\mathit{id}$, $\mathit{course}$, $\mathit{kind}$ (ThreadTarget), $\mathit{targetID}$ (String), $\mathit{author}$ (User), $\mathit{subject}$ (String), $\mathit{content}$ (String, the first post), and $\mathit{replies}$ ($\operatorname{List}\langle\text{Reply}\rangle$).

\subsubsection{Reply}
A record of one reply in a thread, with fields $\mathit{author}$ (User) and $\mathit{content}$ (String). M\_Persistence records when it was written.


\section[MIS of Hardware-Hiding Module]{MIS of Hardware-Hiding Module (\moduleid{M_HardwareHiding})}
\label{M_HardwareHiding}

\wss{You can reference SRS labels, such as R\ref{R_Inputs}.}

\wss{It is also possible to use \LaTeX for hypperlinks to external documents.}

\subsection{Module}

\wss{Short name for the module}
This module provides the platform services used by the system and virtual hardware used by the system, including input/output, file storage, and network connections including HTTPS. This module is provided by the OS, browser, and runtime libraries.

Note: Unlike many other modules listed, this module is entirely externally provided. This will not be explicitly implemented.

\subsection{Uses}
\begin{itemize}
  \item \textbf{Server Operating System:} File system access, networking, and process execution. 
  \item \textbf{Python runtime and standard libraries:} Access to OS services for the backend. 
  \item \textbf{Browser:} Provides access to OS services such as display, user input, and network facilities.
\end{itemize}
No Keystone modules.

\subsection{Syntax}

\subsubsection{Exported Constants}
None

\subsubsection{Exported Access Programs}
None. Provided by the OS, browser, and runtime libraries.

\subsection{Semantics}

\subsubsection{State Variables}

\wss{Not all modules will have state variables.  State variables give the module
  a memory.}
Platform state managed by the OS, browser, and runtime libraries.

\subsubsection{Environment Variables}

\wss{This section is not necessary for all modules.  Its purpose is to capture
  when the module has external interaction with the environment, such as for a
  device driver, screen interface, keyboard, file, etc.}
File system, network connections, user input devices, and browser display.

\subsubsection{Assumptions}

\wss{Try to minimize assumptions and anticipate programmer errors via
  exceptions, but for practical purposes assumptions are sometimes appropriate.}

\begin{itemize}
  \item A supported OS, Python runtime, and browser are available.
  \item Required filesystem and network permissions are configured.
  \item HTTPS is configured for the deployed browser-server communication system. 
\end{itemize}

\subsubsection{Access Routine Semantics}
\wss{A module without environment variables or state variables is unlikely to
  have a state transition.  In this case a state transition can only occur if
  the module is changing the state of another module.}

\wss{Modules rarely have both a transition and an output.  In most cases you
  will have one or the other.}
None

\subsubsection{Local Functions}

\wss{As appropriate} \wss{These functions are for the purpose of specification.
  They are not necessarily something that is going to be implemented
  explicitly.  Even if they are implemented, they are not exported; they only
  have local scope.}
None

\newpage

\section[MIS of User Interface Module]{MIS of User Interface Module (\moduleid{M_UserInterface})}
\label{M_UserInterface}

\subsection{Module}

This module manages Keystone's user interface and views, displaying pages and forms, collecting user input, and presenting results and error messages. It handles presentation but delegates business logic and authorization to the appropriate modules.

\subsection{Uses}

\begin{itemize}
    \item M\_HardwareHiding
    \item M\_Authentication
    \item M\_Authorization
    \item M\_CourseManagement
    \item M\_Enrollment
    \item M\_Projects
    \item M\_StudentsTeams
    \item M\_Recommendation
    \item M\_ApprovalWorkflow
    \item M\_Threads
    \item M\_Notifications
    \item M\_Checklist
    \item M\_CourseReport
\end{itemize}

The UI makes API calls to every backend module listed above to retrieve data and perform the actions available in its views and forms. It displays the results and errors returned by those calls. The access routines below abstract this interaction into presentation and input operations rather than listing each API call separately. The UI uses Hardware Hiding for browser display, input, and networking facilities.

\subsection{Syntax}

\subsubsection{Exported Constants}
None.

\subsubsection{Exported Access Programs}

These proposed interfaces handle presentation and input only. Business operations remain in their owning modules.

\begingroup
\small
\setlength{\tabcolsep}{4pt}
\renewcommand{\arraystretch}{1.2}
\begin{longtable}{>{\raggedright\arraybackslash}p{0.25\textwidth}
                  >{\raggedright\arraybackslash}p{0.32\textwidth}
                  >{\raggedright\arraybackslash}p{0.16\textwidth}
                  >{\raggedright\arraybackslash}p{0.16\textwidth}}
\toprule
\textbf{Name} & \textbf{In} & \textbf{Out} & \textbf{Exceptions} \\
\midrule
\endfirsthead
\toprule
\textbf{Name} & \textbf{In} & \textbf{Out} & \textbf{Exceptions} \\
\midrule
\endhead
\endfoot
\texttt{displayView} & viewName: String, data: ViewData & -- & UnknownView \\
\texttt{displayForm} & formName: String, initialValues: FieldValues & -- & UnknownForm \\
\texttt{readForm} & -- & FieldValues & NoActiveForm \\
\texttt{displayFieldErrors} & errors: FieldErrors & -- & -- \\
\texttt{showMessage} & text: String & -- & -- \\
\texttt{setLoading} & waiting: bool & -- & -- \\
\bottomrule
\caption{Exported access programs of the User Interface module}
\label{TblUIAccessPrograms}
\end{longtable}
\endgroup

\subsection{Semantics}

\subsubsection{State Variables}
\begin{itemize}
  \item $\mathit{currentView}:\text{String}$ - current view which is initially SignIn.
  \item $\mathit{displayedData}:\text{ViewData}$ - values supplied for the view and initially empty.
  \item $\mathit{currentForm}:\text{String}\cup\{\text{None}\}$ - active form which is initially None.
  \item $\mathit{formValues}:\text{FieldValues}$ - map of field names to entered values and the set is initially $\emptyset$.
  \item $\mathit{fieldErrors}:\text{FieldErrors}$ - map of field names to error messages and the set is initially $\emptyset$.
  \item $\mathit{message}:\text{String}$ - feedback which initially $\epsilon$.
  \item $\mathit{loading}:\text{bool}$ - waiting indicator and initially false.
\end{itemize}

\subsubsection{Environment Variables}
Browser display, keyboard, mouse or touch input, file-selection controls, and browser network facilities. The environment variables are external resources instead of application configuration variables.

\subsubsection{Assumptions}
\begin{itemize}
  \item A supported browser is available.
  \item Views cover the tasks assigned to the UI by the MG. The caller supplies data matching the selected view/form schema and the owning module's types.
  \item Server modules authorize displayed data and independently validate submitted inputs and enforce business rules.
\end{itemize}

\subsubsection{Access Routine Semantics}

Equations follow Section~\ref{SecNotation}. Successful transitions are rendered in the browser. All routines have no program output or exception unless stated.

\paragraph{\texttt{displayView(viewName, data)}}
\begin{itemize}
  \item transition:
  \[
  viewName\in\mathcal{V}\Rightarrow
  \left\{
  \begin{aligned}
    \mathit{currentView}' &= viewName, & \mathit{displayedData}' &= data,\\
    \mathit{currentForm}' &= \text{None}, & \mathit{formValues}' &= \emptyset,\\
    \mathit{fieldErrors}' &= \emptyset, & \mathit{message}' &= \epsilon,\\
    \mathit{loading}' &= \text{false}. &&
  \end{aligned}
  \right.
  \]
  \item output: None.
  \item exception:
  $viewName\notin\mathcal{V}\Rightarrow exc=\text{UnknownView}$.
\end{itemize}

\paragraph{\texttt{displayForm(formName, initialValues)}}
\begin{itemize}
  \item transition:
  \[
  formName\in\mathcal{F}\Rightarrow
  \left\{
  \begin{aligned}
    \mathit{currentForm}' &= formName, & \mathit{formValues}' &= initialValues,\\
    \mathit{fieldErrors}' &= \emptyset, & \mathit{message}' &= \epsilon.
  \end{aligned}
  \right.
  \]
  Render the form and mark required fields without submitting values.
  \item output: None.
  \item exception:
  $formName\notin\mathcal{F}\Rightarrow exc=\text{UnknownForm}$.
\end{itemize}

\paragraph{\texttt{readForm()}}
\begin{itemize}
  \item transition: None since values are neither cleared nor submitted.
  \item output:
  $\mathit{currentForm}\ne\text{None}\Rightarrow out=\mathit{formValues}$.
  \item exception:
  $\mathit{currentForm}=\text{None}\Rightarrow exc=\text{NoActiveForm}$.
\end{itemize}

\paragraph{\texttt{displayFieldErrors(errors)}}
\begin{itemize}
  \item transition:
  \[
    \mathit{fieldErrors}'=errors
    \land \mathit{formValues}'=\mathit{formValues}.
  \]
  Display messages beside visible fields and in an accessible summary, including errors for hidden fields.
  \item output: None.
  \item exception: None.
\end{itemize}

\paragraph{\texttt{showMessage(text)}}
\begin{itemize}
  \item transition: $\mathit{message}'=text$. Display accessible feedback and $text=\epsilon\Rightarrow$ clear the message. Exclude credentials, tokens, and restricted details.
  \item output: None.
  \item exception: None.
\end{itemize}

\paragraph{\texttt{setLoading(waiting)}}
\begin{itemize}
  \item transition: $\mathit{loading}'=waiting$, with browser effects:
  \[
  \begin{aligned}
    waiting &\Rightarrow (\text{indicator shown}\land
                          \text{conflicting submissions disabled}),\\
    \neg waiting &\Rightarrow (\text{indicator hidden}\land
                          \text{otherwise permitted controls restored}).
  \end{aligned}
  \]
  \item output: None.
  \item exception: None.
\end{itemize}

\subsubsection{Local Functions}
\begin{description}
  \item[renderCurrentView(): --] Renders the current view and controls from UI state.
  \item[collectFormValues(): FieldValues] Returns the current browser form values and updates $\mathit{formValues}$.
\end{description}

\newpage


\section[MIS of Authentication Module]{MIS of Authentication Module (\moduleid{M_Authentication})}
\label{M_Authentication}

\subsection{Module}

This module authenticates Keystone users through McMaster SSO or a trusted external provider. It maintains user identities and authenticated sessions, verifies credentials, and re-verifies users before high-risk actions. Role-based permissions are handled by M\_Authorization.

\subsection{Uses}

\begin{itemize}
    \item M\_HardwareHiding
    \item M\_Persistence
    \item M\_Authorization
    \item M\_AuditLog
\end{itemize}

\subsection{Syntax}

\subsubsection{Exported Constants}
None.

\subsubsection{Exported Access Programs}

\begingroup
\small
\setlength{\tabcolsep}{4pt}
\renewcommand{\arraystretch}{1.2}
\begin{longtable}{>{\raggedright\arraybackslash}p{0.24\textwidth}
                  >{\raggedright\arraybackslash}p{0.34\textwidth}
                  >{\raggedright\arraybackslash}p{0.16\textwidth}
                  >{\raggedright\arraybackslash}p{0.16\textwidth}}
\toprule
\textbf{Name} & \textbf{In} & \textbf{Out} & \textbf{Exceptions} \\
\midrule
\endfirsthead
\toprule
\textbf{Name} & \textbf{In} & \textbf{Out} & \textbf{Exceptions} \\
\midrule
\endhead
\endfoot
\texttt{signIn} & provider: AuthProvider, proof: AuthProof & SessionID & AuthenticationFailed, ProviderUnavailable \\
\texttt{signOut} & session: SessionID & -- & InvalidSession \\
\texttt{isAuthenticated} & session: SessionID & bool & -- \\
\texttt{getIdentity} & session: SessionID & User & InvalidSession \\
\texttt{reverify} & session: SessionID, proof: AuthProof & bool & InvalidSession, ProviderUnavailable \\
\bottomrule
\caption{Exported access programs of the Authentication module}
\label{TblAuthenticationAccessPrograms}
\end{longtable}
\endgroup

\subsection{Semantics}

\subsubsection{State Variables}

\begin{itemize}
    \item $\mathit{sessions}:\operatorname{Map}\langle\text{SessionID},\text{User}\rangle$ - active authenticated sessions, initially $\emptyset$.

    \item $\mathit{identities}:\operatorname{Map}\langle\text{String},\text{User}\rangle$ - stored user identities indexed by linked provider account, initially loaded from M\_Persistence.
\end{itemize}

\subsubsection{Environment Variables}

McMaster SSO, the external authentication provider, server network connections, and persistent identity storage.

\subsubsection{Assumptions}

\begin{itemize}
    \item Trusted identity providers authenticate users and return verifiable identity information.
    \item All authentication communication uses HTTPS.
\end{itemize}

\subsubsection{Access Routine Semantics}

\paragraph{\texttt{signIn(provider, proof)}}
\begin{itemize}
    \item transition: If the provider successfully verifies the proof, find the linked Keystone user or create one. For a first-time external-provider sign-in, assign the external supervisor role with Unverified status. Generate a unique session identifier $s$:
    \[
    \begin{aligned}
        \mathit{sessions}' &=
        \mathit{sessions}\cup\{s\mapsto user\}.
    \end{aligned}
    \]
    Store new identity and account information through M\_Persistence. The new external supervisor has only the unverified access scope in the SRS.

    \item output: $out=s$.

    \item exception:
    \begin{itemize}
        \item AuthenticationFailed if identity verification fails or account creation fails.
        \item ProviderUnavailable if the provider cannot be reached.
    \end{itemize}
\end{itemize}

\paragraph{\texttt{signOut(session)}}
\begin{itemize}
    \item transition:
    \[
    \mathit{sessions}'=
        \mathit{sessions}\setminus\{session\mapsto\mathit{sessions}[session]\}.
    \]
    \item output: None.
    \item exception:
    $session\notin\operatorname{dom}(\mathit{sessions})\Rightarrow exc=\text{InvalidSession}$.
\end{itemize}

\paragraph{\texttt{isAuthenticated(session)}}
\begin{itemize}
    \item transition: None.
    \item output:
    \[
    out=
    \begin{cases}
        \text{true}, & \operatorname{validSession}(session),\\
        \text{false}, & \text{otherwise}.
    \end{cases}
    \]
    \item exception: None.
\end{itemize}

\paragraph{\texttt{getIdentity(session)}}
\begin{itemize}
    \item transition: None.
    \item output:
    \[
    \operatorname{validSession}(session)
    \Rightarrow out=\mathit{sessions}[session].
    \]
    \item exception:
    \[
    \neg\operatorname{validSession}(session)
    \Rightarrow exc=\text{InvalidSession}.
    \]
\end{itemize}

\paragraph{\texttt{reverify(session, proof)}}
\begin{itemize}
    \item transition: None.
    \item output:
    \[
    out=
    \operatorname{verifyProof}
    (\mathit{sessions}[session],proof).
    \]
    The requesting module calls re-verification for the triggers listed in \texttt{FR\_ReverifyUsers}. The proof uses re-entered credentials, an email code, an MFA challenge, or an SSO challenge.

    \item exception:
    \begin{itemize}
    \item InvalidSession if the session is invalid.
    \item ProviderUnavailable if external verification cannot be completed.
    \end{itemize}
\end{itemize}

\subsubsection{Local Functions}

\begin{description}
    \item[validSession(session: SessionID): bool] Returns true if the session exists and remains valid.
    \item[verifyProof(user: User, proof: AuthProof): bool] Returns true if the proof verifies for the specified user.
    \item[hashPassword(password: String): String] Returns a salted Argon2id hash of the password.
    \item[verifyPassword(password: String, hash: String): bool] Returns true if the password matches the stored hash without storing the plaintext password.
\end{description}

\newpage

\section[MIS of Authorization Module]{MIS of Authorization Module (\moduleid{M_Authorization})}
\label{M_Authorization}

\subsection{Module}

This module manages user roles and access permissions. It determines whether authenticated users may perform actions or access records based on their roles, course permissions, and external supervisor verification status. It also handles authorized changes to user roles and verification status.

\subsection{Uses}

\begin{itemize}
    \item M\_Persistence
    \item M\_AuditLog
\end{itemize}

\subsection{Syntax}

\subsubsection{Exported Constants}
None.

\subsubsection{Exported Access Programs}

\begingroup
\small
\setlength{\tabcolsep}{4pt}
\renewcommand{\arraystretch}{1.2}
\begin{longtable}{>{\raggedright\arraybackslash}p{0.28\textwidth}
                  >{\raggedright\arraybackslash}p{0.34\textwidth}
                  >{\raggedright\arraybackslash}p{0.12\textwidth}
                  >{\raggedright\arraybackslash}p{0.16\textwidth}}
\toprule
\textbf{Name} & \textbf{In} & \textbf{Out} & \textbf{Exceptions} \\
\midrule
\endfirsthead
\toprule
\textbf{Name} & \textbf{In} & \textbf{Out} & \textbf{Exceptions} \\
\midrule
\endhead
\endfoot
\texttt{hasAccess} & user: User, course: Course, action: Action, record: Value & bool & -- \\
\texttt{getRole} & user: User, course: Course & UserRole & RoleNotFound \\
\texttt{updateRole} & actor: User, user: User, course: Course, role: UserRole & -- & Unauthorized \\
\texttt{verifyExternalSupervisor} & user: User, project: Project & -- & InvalidVerification \\
\bottomrule
\caption{Exported access programs of the Authorization module}
\label{TblAuthorizationAccessPrograms}
\end{longtable}
\endgroup

\subsection{Semantics}

\subsubsection{State Variables}

\begin{itemize}
    \item $\mathit{roles}:
    \operatorname{Map}\langle
    \text{User},
    \operatorname{Map}\langle\text{Course},\text{UserRole}\rangle
    \rangle$
    - assigned course roles, initially loaded from
    M\_Persistence.

    \item $\mathit{verification}:
    \operatorname{Map}\langle
    \text{User},\text{SupervisorVerification}\rangle$
    - external supervisor verification states, initially
    loaded from M\_Persistence.

    \item $\mathit{administrators}:
    \operatorname{Map}\langle\text{User},\text{bool}\rangle$
    - the SRS \texttt{is\_admin} permission flag for Instructor and TA accounts. It is not a separate role and is initially loaded from M\_Persistence.

    \item $\mathit{permittedActions}:\operatorname{Map}\langle\text{UserRole},\operatorname{Set}\langle\text{Action}\rangle\rangle$ - actions allowed for each role, initially loaded from M\_Persistence.
\end{itemize}

\subsubsection{Environment Variables}

Persistent user roles, course associations, supervisor verification records, and the \texttt{is\_admin} permission flags for Instructor and TA accounts.

\subsubsection{Assumptions}

\begin{itemize}
    \item User identities are authenticated through M\_Authentication before authorization checks.

    \item Access decisions follow the role permissions and visibility requirements defined in the SRS.

    \item The requesting module authenticates the actor and completes any required re-verification before requesting a role change.
\end{itemize}

\subsubsection{Access Routine Semantics}

\paragraph{\texttt{hasAccess(user, course, action, record)}}
\begin{itemize}
    \item transition: None.

    \item output: $out=(course\in\operatorname{dom}(\mathit{roles}[user]))\land(action\in\mathit{permittedActions}[\mathit{roles}[user][course]])\land\operatorname{recordInScope}(user,course,action,record)$. If the user has no course role, the result is false.

    \item exception: None.
\end{itemize}

\paragraph{Record scope}
$\operatorname{recordInScope}(user:User,course:Course,action:Action,record:Value):\text{bool}$ returns true when the record belongs to the course and the action is allowed for that record under the SRS permissions and visibility rules (see SRS tables \srsref{TAB-Visibility} and \srsref{TAB-ExternalAccess}). For reads, an unverified external supervisor may access course reference material, proposal templates, and their own proposals, projects, and threads. A verified external supervisor may also access applications to their own projects with applicant team profiles, student and team profiles, and student-proposed projects. Neither may access another supervisor's proposals or projects, which must not be revealed in responses, catalogs, searches, counts, statuses, or notifications.

\paragraph{\texttt{getRole(user, course)}}
\begin{itemize}
    \item transition: None.

    \item output:
    \[
    course\in\operatorname{dom}(\mathit{roles}[user])
    \Rightarrow
    out=\mathit{roles}[user][course].
    \]

    \item exception: RoleNotFound if the user has
    no assigned role in the specified course.
\end{itemize}

\paragraph{\texttt{updateRole(actor, user, course, role)}}
\begin{itemize}
    \item transition: If the requesting module has authenticated and re-verified the actor as required by \texttt{FR\_ReverifyUsers}, and
    \[
    \operatorname{hasAccess}(actor,course,\text{Update},user)=\text{true},
    \]
    \[
    \mathit{roles}' =
    \mathit{roles}\oplus
    \left\{
    user\mapsto
    \left(
    \mathit{roles}[user]\oplus
    \{course\mapsto role\}
    \right)
    \right\}.
    \]
    If the user's current role in the course is Supervisor and the new role is not, set
    \[
    \mathit{verification}'=\mathit{verification}\oplus\{user\mapsto\text{Unverified}\}.
    \]
    Otherwise, $\mathit{verification}'=\mathit{verification}$.
    Changing a verified external supervisor's role away from Supervisor resets their verification state to Unverified, reversing their verified supervisor access. Persist and audit both changes.

    \item output: None.

    \item exception: Unauthorized if the actor is not authorized or the required authentication or re-verification has not succeeded.
\end{itemize}

\paragraph{\texttt{verifyExternalSupervisor(user, project)}}
\begin{itemize}
    \item transition: If the user is an external supervisor and an instructor approves a project proposal submitted by that user:
    \[
    \mathit{verification}' =
    \mathit{verification}\oplus
    \{user\mapsto\text{Verified}\}.
    \]
    The change is persisted and audited.

    Approval of a proposal submitted by another user, a request for changes, or rejection does not verify the supervisor. Withdrawal of an approved proposal does not reverse verification.

    \item output: None.

    \item exception: InvalidVerification if the required approval conditions are not satisfied.
\end{itemize}

\subsubsection{Local Functions}
None.

\newpage

\section[MIS of Course Management Module]{MIS of Course Management Module (\moduleid{M_CourseManagement})}
\label{M_CourseManagement}

\subsection{Module}

This module manages course-specific settings, enabled
features, and reference materials. It allows instructors
to configure course functionality, prevents features
from being disabled when outstanding work exists,
and allows instructors and TAs to maintain and publish
course resources for students and supervisors.

\subsection{Uses}

\begin{itemize}
    \item M\_Authorization
    \item M\_Authentication
    \item M\_Persistence
    \item M\_AuditLog
\end{itemize}

\subsection{Syntax}

\subsubsection{Exported Constants}
None.

\subsubsection{Exported Access Programs}

\begingroup
\small
\setlength{\tabcolsep}{4pt}
\renewcommand{\arraystretch}{1.2}
\begin{longtable}{>{\raggedright\arraybackslash}p{0.24\textwidth}
                  >{\raggedright\arraybackslash}p{0.34\textwidth}
                  >{\raggedright\arraybackslash}p{0.16\textwidth}
                  >{\raggedright\arraybackslash}p{0.16\textwidth}}
\toprule
\textbf{Name} & \textbf{In} & \textbf{Out} & \textbf{Exceptions} \\
\midrule
\endfirsthead
\toprule
\textbf{Name} & \textbf{In} & \textbf{Out} & \textbf{Exceptions} \\
\midrule
\endhead
\endfoot
\texttt{getSettings} & user: User, course: Course & Dictionary & Unauthorized \\
\texttt{updateSetting} & user: User, course: Course, key: String, value: Value, session: SessionID, proof: AuthProof & -- & Unauthorized, InvalidSetting \\
\texttt{isFeatureEnabled} & user: User, course: Course, feature: Feature & bool & Unauthorized \\
\texttt{setFeatureEnabled} & user: User, course: Course, feature: Feature, enabled: bool & -- & Unauthorized, PendingWork \\
\texttt{getPendingWork} & user: User, course: Course, feature: Feature & List & Unauthorized \\
\texttt{recordPendingWork} & course: Course, feature: Feature, recordID: String, pending: bool & -- & -- \\
\texttt{saveResource} & user: User, course: Course, resource: CourseResource & -- & Unauthorized, InvalidResource \\
\texttt{publishResource} & user: User, course: Course, resourceID: String & -- & Unauthorized, ResourceNotFound \\
\texttt{getResources} & user: User, course: Course & $\operatorname{List}\langle\text{CourseResource}\rangle$ & Unauthorized \\
\bottomrule
\caption{Exported access programs of the Course Management module}
\label{TblCourseManagementAccessPrograms}
\end{longtable}
\endgroup

\subsection{Semantics}

\subsubsection{State Variables}

\begin{itemize}
    \item $\mathit{settings}:
    \operatorname{Map}\langle
    \text{Course},\text{Dictionary}\rangle$
    - course-specific configuration values, initially
    loaded from M\_Persistence.

    \item $\mathit{features}:
    \operatorname{Map}\langle
    \text{Course},
    \operatorname{Map}\langle\text{Feature},\text{bool}\rangle
    \rangle$
    - enabled states of supported features for each
    course, initially loaded from M\_Persistence.

    \item $\mathit{pendingWork}:
    \operatorname{Map}\langle
    \text{Course},
    \operatorname{Map}\langle
            \text{Feature},\operatorname{Set}\langle\text{String}\rangle
    \rangle\rangle$
    - identifiers of outstanding records associated
    with each feature, initially loaded from
    M\_Persistence.

    \item $\mathit{resources}:
    \operatorname{Map}\langle
    \text{String},\text{CourseResource}\rangle$
    - stored course resources indexed by their
    identifiers, initially loaded from M\_Persistence.
\end{itemize}

\subsubsection{Environment Variables}

Persistent course configuration, feature states,
outstanding work records, and course resource storage.

\subsubsection{Assumptions}

\begin{itemize}
    \item M\_Authorization determines user permissions for course configuration and resource management.
    \item Changes to course rules and approval modes require successful re-verification through M\_Authentication.
\end{itemize}

\subsubsection{Access Routine Semantics}

\paragraph{\texttt{getSettings(user, course)}}
\begin{itemize}
    \item transition: None.

    \item output: If $\operatorname{hasAccess}(user,course,\text{Read},course)=\text{true}$:
    \[
    out=\mathit{settings}[course].
    \]

    \item exception: Unauthorized if $\operatorname{hasAccess}(user,course,\text{Read},course)=\text{false}$.
\end{itemize}

\paragraph{\texttt{updateSetting(user, course, key, value, session, proof)}}
\begin{itemize}
    \item transition: If
    \[
    \begin{aligned}
    &\operatorname{hasAccess}(user,course,\text{Update},course)=\text{true}
    \land validSetting(course,key,value)\\
    &\land (\operatorname{requiresReverification}(key)\Rightarrow
    \operatorname{reverify}(session,proof)=\text{true}),
    \end{aligned}
    \]
    \[
    \mathit{settings}'=
    \mathit{settings}\oplus
    \left\{
    course\mapsto
    \left(
    \mathit{settings}[course]\oplus
    \{key\mapsto value\}
    \right)
    \right\}.
    \]
    Persist and audit the change.

    \item output: None.

    \item exception:
    \begin{itemize}
        \item Unauthorized if $\operatorname{hasAccess}(user,course,\text{Update},course)=\text{false}$ or the required re-verification fails.
        \item InvalidSetting if the setting is unsupported or the value is invalid.
    \end{itemize}
\end{itemize}

\paragraph{\texttt{isFeatureEnabled(user, course, feature)}}
\begin{itemize}
    \item transition: None.

    \item output: If $\operatorname{hasAccess}(user,course,\text{Read},course)=\text{true}$:
    \[
    out=\mathit{features}[course][feature].
    \]

    \item exception: Unauthorized if $\operatorname{hasAccess}(user,course,\text{Read},course)=\text{false}$.
\end{itemize}

\paragraph{\texttt{setFeatureEnabled(user, course, feature, enabled)}}
\begin{itemize}
    \item transition: If $\operatorname{hasAccess}(user,course,\text{Configure},course)=\text{true}$ and disabling the feature will not strand outstanding work:
    \[
    \mathit{features}'=
    \mathit{features}\oplus
    \left\{
    course\mapsto
    \left(
    \mathit{features}[course]\oplus
    \{feature\mapsto enabled\}
    \right)
    \right\}.
    \]
    Persist and audit the change.

    \item output: None.

    \item exception:
    \begin{itemize}
        \item Unauthorized if $\operatorname{hasAccess}(user,course,\text{Configure},course)=\text{false}$.
        \item PendingWork if
        \[
        \neg enabled\land
        \mathit{pendingWork}[course][feature]
        \ne\emptyset.
        \]
    \end{itemize}
    When PendingWork occurs, no feature state changes. The caller may use getPendingWork to display the outstanding records.
\end{itemize}

\paragraph{\texttt{getPendingWork(user, course, feature)}}
\begin{itemize}
    \item transition: None.

    \item output: If $\operatorname{hasAccess}(user,course,\text{Read},course)=\text{true}$,
    return the outstanding record identifiers:
    \[
    out=\operatorname{List}
    (\mathit{pendingWork}[course][feature]).
    \]

    \item exception: Unauthorized if $\operatorname{hasAccess}(user,course,\text{Read},course)=\text{false}$.
\end{itemize}

\paragraph{\texttt{recordPendingWork(course, feature, recordID, pending)}}
\begin{itemize}
    \item transition: Called only by a trusted Keystone workflow module. Let
    \[
    W=\mathit{pendingWork}[course][feature].
    \]
    Then:
    \[
    W'=
    \begin{cases}
        W\cup\{recordID\}, & pending,\\
        W\setminus\{recordID\}, & \neg pending.
    \end{cases}
    \]
    Replace the corresponding pending-work entry with $W'$ and persist the change.

    \item output: None.

    \item exception: None.
\end{itemize}

\paragraph{\texttt{saveResource(user, course, resource)}}
\begin{itemize}
    \item transition: If $\operatorname{hasAccess}(user,course,\text{Create},course)=\text{true}$ and the resource is valid:
    \[
    \mathit{resources}'=
    \mathit{resources}\oplus
    \{resource.id\mapsto resource\}.
    \]
    The resource is associated with the supplied course. New resources remain unpublished until explicitly published. Persist and audit the change.

    \item output: None.

    \item exception:
    \begin{itemize}
        \item Unauthorized if $\operatorname{hasAccess}(user,course,\text{Create},course)=\text{false}$.
        \item InvalidResource if required resource information is invalid.
    \end{itemize}
\end{itemize}

\paragraph{\texttt{publishResource(user, course, resourceID)}}
\begin{itemize}
    \item transition: If $\operatorname{hasAccess}(user,course,\text{Update},resourceID)=\text{true}$ and the resource belongs to the course:
    \[
    \mathit{resources}'[resourceID].published
    =\text{true}.
    \]
    Other resource fields remain unchanged. Persist and audit the change.

    \item output: None.

    \item exception:
    \begin{itemize}
        \item Unauthorized if $\operatorname{hasAccess}(user,course,\text{Update},resourceID)=\text{false}$.
        \item ResourceNotFound if the resource does not exist in the course.
    \end{itemize}
\end{itemize}

\paragraph{\texttt{getResources(user, course)}}
\begin{itemize}
    \item transition: None.

    \item output: If $\operatorname{hasAccess}(user,course,\text{Read},course)=\text{true}$, return the list of published resources belonging to the requested course:
    \[
    out=[r\mid r\in\mathit{resources},
    \ r.course=course\land r.published].
    \]
    Unpublished resources are not returned to students or supervisors.

    \item exception: Unauthorized if $\operatorname{hasAccess}(user,course,\text{Read},course)=\text{false}$.
\end{itemize}

\subsubsection{Local Functions}

\begin{description}
    \item[validSetting(course: Course, key: String, value: Value): bool] Returns true if the setting exists and the supplied value satisfies its constraints.
    \item[requiresReverification(key: String): bool] Returns true when the setting changes course rules or the approval mode, as required by \texttt{FR\_ReverifyUsers}.

    \item[validResource(resource: CourseResource): bool] Returns true if the resource contains the required information and valid content.
\end{description}


\newpage

\section[MIS of Enrollment Module]{MIS of Enrollment Module (\moduleid{M_Enrollment})}
\label{M_Enrollment}

\subsection{Module}

This module manages course enrollment records. It imports student lists from Avenue to Learn or uploaded files, validates and previews enrollment changes, and applies imports only after confirmation. It also supports manual enrollment updates and flags students missing from the latest imported list without automatically removing them.

\subsection{Uses}

\begin{itemize}
    \item M\_StudentsTeams
    \item M\_CourseManagement
    \item M\_HardwareHiding
    \item M\_AuditLog
    \item M\_Persistence
    \item M\_Authorization
    \item M\_Authentication
\end{itemize}

\subsection{Syntax}

\subsubsection{Exported Constants}
None.

\subsubsection{Exported Access Programs}

\begingroup
\small
\setlength{\tabcolsep}{4pt}
\renewcommand{\arraystretch}{1.2}
\begin{longtable}{>{\raggedright\arraybackslash}p{0.24\textwidth}
                  >{\raggedright\arraybackslash}p{0.34\textwidth}
                  >{\raggedright\arraybackslash}p{0.16\textwidth}
                  >{\raggedright\arraybackslash}p{0.16\textwidth}}
\toprule
\textbf{Name} & \textbf{In} & \textbf{Out} & \textbf{Exceptions} \\
\midrule
\endfirsthead
\toprule
\textbf{Name} & \textbf{In} & \textbf{Out} & \textbf{Exceptions} \\
\midrule
\endhead
\endfoot
\texttt{previewImport} & user: User, course: Course, source: EnrollmentSource, entries: List & EnrollmentPreview & Unauthorized, InvalidSource \\
\texttt{confirmImport} & user: User, course: Course, importID: ImportID, session: SessionID, proof: AuthProof & -- & Unauthorized, InvalidImport, ValidationFailed \\
\texttt{addStudent} & user: User, course: Course, entry: EnrollmentEntry, session: SessionID, proof: AuthProof & -- & Unauthorized, InvalidEntry, DuplicateStudent \\
\texttt{getEnrollment} & user: User, course: Course & List & Unauthorized \\
\texttt{getDiscrepancies} & user: User, course: Course & List & Unauthorized \\
\texttt{confirmRemoval} & user: User, course: Course, student: User, session: SessionID, proof: AuthProof & -- & Unauthorized, StudentNotFlagged \\
\bottomrule
\caption{Exported access programs of the Enrollment module}
\label{TblEnrollmentAccessPrograms}
\end{longtable}
\endgroup

\subsection{Semantics}

\subsubsection{State Variables}

\begin{itemize}
    \item $\mathit{enrollment}:
    \operatorname{Map}\langle
    \text{Course},
    \operatorname{Map}\langle
    \text{User},\text{EnrollmentStatus}
    \rangle\rangle$
    - current enrollment records and statuses,
    initially loaded from M\_Persistence.

    \item $\mathit{pendingImports}:
    \operatorname{Map}\langle
    \text{ImportID},\text{EnrollmentPreview}\rangle$
    - validated previews awaiting confirmation,
    initially $\emptyset$.
\end{itemize}

\subsubsection{Environment Variables}

Avenue to Learn, uploaded enrollment files, persistent enrollment records, and server network connections.

\subsubsection{Assumptions}

\begin{itemize}
    \item Enrollment data contains only the fields permitted by DPR\_MinimalEnrollmentData.
    \item Enrollment calls M\_Authentication.reverify before an enrollment import is confirmed.
\end{itemize}

\subsubsection{Access Routine Semantics}

\paragraph{\texttt{previewImport(user, course, source, entries)}}
\begin{itemize}
    \item transition: If $\operatorname{hasAccess}(user,course,\text{Import},course)=\text{true}$, retrieve the enrollment list using the specified source and compare its entries with the existing enrollment.

    For a valid import, generate a unique identifier
    $i$ and preview $p$:
    \[
    \mathit{pendingImports}' =
    \mathit{pendingImports}
    \cup\{i\mapsto p\}.
    \]

    The preview identifies new students, matching students, students missing from the imported list, duplicate identifiers, and malformed entries.

    No enrollment records are changed.

    \item output: $out=p$.

    \item exception:
    \begin{itemize}
        \item Unauthorized if $\operatorname{hasAccess}(user,course,\text{Import},course)=\text{false}$.
        \item InvalidSource if the requested source cannot provide a readable enrollment list.
    \end{itemize}
\end{itemize}

\paragraph{\texttt{confirmImport(user, course, importID, session, proof)}}
\begin{itemize}
    \item transition: If
    \[
    \begin{aligned}
    &\operatorname{hasAccess}(user,course,\text{ConfirmImport},importID)=\text{true}
    \land \operatorname{reverify}(session,proof)=\text{true}\\
    &\land importID\in\operatorname{dom}(\mathit{pendingImports})
    \land \mathit{pendingImports}[importID].course=course,
    \end{aligned}
    \]
    apply the validated import.

    Let $I$ be the set of students in the import
    and $E$ the set of currently enrolled students.

    For every $u\in I$:
    \[
    \mathit{enrollment}'[course][u]
    =\text{Enrolled}.
    \]

    For every $u\in E\setminus I$:
    \[
    \mathit{enrollment}'[course][u]
    =\text{Flagged}.
    \]

    Students absent from the import remain associated with the course and retain their team memberships.

    The changes are saved atomically, audited, and the completed import is removed from pendingImports.

    \item output: None.

    \item exception:
    \begin{itemize}
        \item Unauthorized if $\operatorname{hasAccess}(user,course,\text{ConfirmImport},importID)=\text{false}$ or $\operatorname{reverify}(session,proof)=\text{false}$.
        \item InvalidImport if the preview does not exist, belongs to another course, or is no longer current.
        \item ValidationFailed if the import contains unresolved duplicate or malformed identifiers.
    \end{itemize}
\end{itemize}

\paragraph{\texttt{addStudent(user, course, entry, session, proof)}}
\begin{itemize}
    \item transition: If $\operatorname{hasAccess}(user,course,\text{Create},course)=\text{true}\land\operatorname{reverify}(session,proof)=\text{true}$ and the entry is valid, add the student to the specified course.

    For the identified student $s$:
    \[
    \mathit{enrollment}'[course]
    =
    \mathit{enrollment}[course]
    \oplus\{s\mapsto\text{Enrolled}\}.
    \]

    A previously flagged student may be
    restored to Enrolled status.
    Persist and audit the change.

    \item output: None.

    \item exception:
    \begin{itemize}
        \item Unauthorized if $\operatorname{hasAccess}(user,course,\text{Create},course)=\text{false}$ or $\operatorname{reverify}(session,proof)=\text{false}$.
        \item InvalidEntry if the supplied enrollment information is malformed or incomplete.
        \item DuplicateStudent if the student is already enrolled in the course.
    \end{itemize}
\end{itemize}

\paragraph{\texttt{getEnrollment(user, course)}}
\begin{itemize}
    \item transition: None.

    \item output: If $\operatorname{hasAccess}(user,course,\text{Read},course)=\text{true}$, return the course's
    student enrollment records and statuses:
    \[
    out=\operatorname{List}
    (\mathit{enrollment}[course]).
    \]

    \item exception: Unauthorized if $\operatorname{hasAccess}(user,course,\text{Read},course)=\text{false}$.
\end{itemize}

\paragraph{\texttt{getDiscrepancies(user, course)}}
\begin{itemize}
    \item transition: None.

    \item output: If $\operatorname{hasAccess}(user,course,\text{Read},course)=\text{true}$, return students requiring review:
    \[
    out=[
    u\mid
    u\in\operatorname{dom}(\mathit{enrollment}[course])
    \land
    \mathit{enrollment}[course][u]
    =\text{Flagged}
    ].
    \]

    \item exception: Unauthorized if $\operatorname{hasAccess}(user,course,\text{Read},course)=\text{false}$.
\end{itemize}

\paragraph{\texttt{confirmRemoval(user, course, student, session, proof)}}
\begin{itemize}
    \item transition: If $\operatorname{hasAccess}(user,course,\text{Delete},student)=\text{true}\land\operatorname{reverify}(session,proof)=\text{true}$ and the student has Flagged status, confirm the discrepancy and initiate the authorized removal of the student's course association.

    The operation must coordinate with modules managing course roles and team memberships. No team membership is silently removed.

    Persist and audit the confirmed change.

    \item output: None.

    \item exception:
    \begin{itemize}
        \item Unauthorized if $\operatorname{hasAccess}(user,course,\text{Delete},student)=\text{false}$ or $\operatorname{reverify}(session,proof)=\text{false}$.
        \item StudentNotFlagged if the student has no flagged enrollment discrepancy in the specified course.
    \end{itemize}
\end{itemize}

\subsubsection{Local Functions}

\begin{description}
    \item[validEntry(entry: EnrollmentEntry): bool]
    Returns true if the entry contains valid student identification, program, and course information.

    \item[findDuplicates(entries: List): List]
    Returns entries with duplicate student identifiers.

    \item[compareEnrollment(course: Course, entries: List): EnrollmentPreview]
    Compares the imported list with current enrollment and produces the proposed additions, matches, missing students, and validation errors.
\end{description}

\newpage

\section[MIS of Projects Module]{MIS of Projects Module (\moduleid{M_Projects})}
\label{M_Projects}

\subsection{Module}

This module manages project records and their life cycle (Figure~\ref{fig:lifecycle} in the SRS). It provides each course's proposal templates, records proposals and their review, tracks how many teams each project can take, and records team applications, supervisor decisions, and agreements. It reserves and releases project slots, withdraws a team's competing applications when the team reaches agreement, and returns the projects a user may see, searched, filtered, and sorted as requested.

\subsection{Uses}

\begin{itemize}
    \item M\_Authorization
    \item M\_CourseManagement
    \item M\_StudentsTeams
    \item M\_Notifications
    \item M\_Persistence
    \item M\_AuditLog
\end{itemize}

\subsection{Syntax}

\subsubsection{Exported Constants}
None.

\subsubsection{Exported Access Programs}

\begingroup
\small
\setlength{\tabcolsep}{4pt}
\renewcommand{\arraystretch}{1.2}
\begin{longtable}{>{\raggedright\arraybackslash}p{0.24\textwidth}
                  >{\raggedright\arraybackslash}p{0.34\textwidth}
                  >{\raggedright\arraybackslash}p{0.16\textwidth}
                  >{\raggedright\arraybackslash}p{0.16\textwidth}}
\toprule
\textbf{Name} & \textbf{In} & \textbf{Out} & \textbf{Exceptions} \\
\midrule
\endfirsthead
\toprule
\textbf{Name} & \textbf{In} & \textbf{Out} & \textbf{Exceptions} \\
\midrule
\endhead
\endfoot
\texttt{getProposalTemplate} & user: User, course: Course, kind: String & $\operatorname{Map}\langle\text{String},\text{bool}\rangle$ & Unauthorized, NoTemplate \\
\texttt{setProposalTemplate} & user: User, course: Course, kind: String, template: $\operatorname{Map}\langle\text{String},\text{bool}\rangle$ & -- & Unauthorized, InvalidTemplate \\
\texttt{createProposal} & user: User, course: Course, values: FieldValues, capacity: $\mathbb{N}$, supervisor: $\text{User}\cup\{\text{None}\}$ & Project & Unauthorized, FeatureDisabled, NoTemplate, NotInTeam, InvalidCapacity, InvalidProposal \\
\texttt{updateProposal} & user: User, project: Project, values: FieldValues & -- & Unauthorized, InvalidStatus, InvalidProposal \\
\texttt{submitProposal} & user: User, project: Project & -- & Unauthorized, FeatureDisabled, InvalidStatus, InvalidProposal \\
\texttt{withdrawProposal} & user: User, project: Project & -- & Unauthorized, InvalidStatus \\
\texttt{deleteProject} & user: User, project: Project & -- & Unauthorized, HasReservedTeams \\
\texttt{reviewProposal} & user: User, project: Project, decision: Action, comment: String & -- & Unauthorized, InvalidStatus, InvalidDecision \\
\texttt{setCapacity} & user: User, project: Project, capacity: $\mathbb{N}$ & -- & Unauthorized, InvalidStatus, InvalidCapacity \\
\texttt{closeSelection} & user: User, course: Course & -- & Unauthorized \\
\texttt{apply} & user: User, project: Project & ProjectApplication & Unauthorized, InvalidStatus, NotInTeam, DuplicateApplication, AlreadyAgreed \\
\texttt{decideApplication} & user: User, application: ProjectApplication, decision: Action, comment: String & -- & Unauthorized, InvalidStatus, InvalidDecision, NoSlotAvailable \\
\texttt{confirmAgreement} & user: User, application: ProjectApplication & -- & Unauthorized, InvalidStatus \\
\texttt{withdrawApplication} & user: User, application: ProjectApplication & -- & Unauthorized, InvalidStatus \\
\texttt{updateApplicationStatus} & application: ProjectApplication, status: TeamProjectStatus & -- & InvalidTransition \\
\texttt{getProject} & user: User, project: Project & Project & Unauthorized \\
\texttt{getApplication} & user: User, application: ProjectApplication & ProjectApplication & Unauthorized \\
\texttt{listProjects} & user: User, course: Course, text: String, statuses: $\operatorname{Set}\langle\text{ProjectStatus}\rangle$, sortBy: String & $\operatorname{List}\langle\text{Project}\rangle$ & Unauthorized, InvalidQuery \\
\texttt{getApplications} & user: User, project: Project & $\operatorname{List}\langle\text{ProjectApplication}\rangle$ & Unauthorized \\
\texttt{getTeamApplications} & user: User, team: Team & $\operatorname{List}\langle\text{ProjectApplication}\rangle$ & Unauthorized \\
\bottomrule
\caption{Exported access programs of the Projects module}
\label{TblProjectsAccessPrograms}
\end{longtable}
\endgroup

\subsection{Semantics}

\subsubsection{State Variables}

\begin{itemize}
    \item $\mathit{templates}:\operatorname{Map}\langle\text{Course},\operatorname{Map}\langle\text{String},\operatorname{Map}\langle\text{String},\text{bool}\rangle\rangle\rangle$ - for each course and each template kind (\texttt{Supervisor} or \texttt{Student}), the proposal fields, each mapped to true if it is required. Initially loaded from M\_Persistence.

    \item $\mathit{projects}:\operatorname{Map}\langle\text{String},\text{Project}\rangle$ - projects indexed by identifier, initially loaded from M\_Persistence.

    \item $\mathit{applications}:\operatorname{Map}\langle\text{String},\text{ProjectApplication}\rangle$ - team applications indexed by identifier, initially loaded from M\_Persistence.
\end{itemize}

\subsubsection{Environment Variables}

Persistent project, application, and template records.

\subsubsection{Assumptions}

\begin{itemize}
    \item Users are authenticated through M\_Authentication before any routine is called.
    \item M\_Authorization enforces the SRS visibility rules through \texttt{hasAccess}, including that no user can see another supervisor's proposals or projects unless the SRS allows it.
    \item The caller has re-verified the user, as required by \texttt{FR\_ReverifyUsers}, before calling \texttt{deleteProject}, \texttt{decideApplication} with Accept, and \texttt{confirmAgreement}.
    \item Every proposal template contains a required field named \texttt{title}.
    \item The records passed in are current, M\_Persistence assigns each new record a unique identifier, and it keeps the \texttt{updated\_at} value of each record.
    \item $\operatorname{notify}(event,record)$ is the M\_Notifications routine that decides the recipients, stores the notification in the caller's transaction, and delivers it later.
    \item \texttt{updateApplicationStatus} is called only by M\_ApprovalWorkflow, and its changes are saved in the caller's transaction. The statuses Applied, SupervisorAccepts, Agreed, and DeclinedOrWithdrawn are set by this module's own routines.
    \item The proposer of a student-proposed project is a team member. How the proposing team is linked to its project, and who names its supervisor, is not decided by the SRS, so \texttt{createProposal} records them in $\mathit{project}.team$ and $\mathit{project}.supervisor$ and nothing more.
\end{itemize}

\subsubsection{Access Routine Semantics}

\paragraph{\texttt{getProposalTemplate(user, course, kind)}}
\begin{itemize}
    \item transition: None.
    \item output: If $\operatorname{hasAccess}(user,course,\text{Read},course)=\text{true}$ and $kind\in\operatorname{dom}(\mathit{templates}[course])$:
    \[
    out=\mathit{templates}[course][kind].
    \]
    \item exception:
    \begin{itemize}
        \item Unauthorized if $\operatorname{hasAccess}(user,course,\text{Read},course)=\text{false}$.
        \item NoTemplate if $kind\notin\operatorname{dom}(\mathit{templates}[course])$.
    \end{itemize}
\end{itemize}

\paragraph{\texttt{setProposalTemplate(user, course, kind, template)}}
\begin{itemize}
    \item transition: If $\operatorname{hasAccess}(user,course,\text{Configure},course)=\text{true}$, $kind\in\{\text{Supervisor},\text{Student}\}$, and $template[\text{title}]=\text{true}$:
    \[
    \mathit{templates}'=\mathit{templates}\oplus\{course\mapsto(\mathit{templates}[course]\oplus\{kind\mapsto template\})\},
    \]
    where $\mathit{templates}[course]=\emptyset$ if the course has no template yet. Existing proposals are not changed. Persist and audit the change.
    \item output: None.
    \item exception:
    \begin{itemize}
        \item Unauthorized if $\operatorname{hasAccess}(user,course,\text{Configure},course)=\text{false}$.
        \item InvalidTemplate if $kind$ is not Supervisor or Student, or \texttt{title} is missing from the template or not required.
    \end{itemize}
\end{itemize}

\paragraph{\texttt{createProposal(user, course, values, capacity, supervisor)}}
\begin{itemize}
    \item transition: Let $r=\operatorname{getRole}(user,course)$ and $k=\operatorname{templateKind}(r)$. If
    \[
    \begin{aligned}
    &\operatorname{hasAccess}(user,course,\text{Create},course)=\text{true}
    \land\operatorname{isFeatureEnabled}(user,course,\text{ProjectProposals})=\text{true}\\
    &\land k\in\operatorname{dom}(\mathit{templates}[course])
    \land\operatorname{dom}(values)\subseteq\operatorname{dom}(\mathit{templates}[course][k])\\
    &\land capacity\ge 1
    \land(r=\text{Student}\Rightarrow\operatorname{getTeamOf}(user,course)\ne\text{None}\land capacity=1),
    \end{aligned}
    \]
    create a project $p$ with $p.course=course$, $p.proposer=user$, $p.fields=values$, $p.capacity=capacity$, $p.reserved=\emptyset$, and $p.status=\text{Draft}$. If $r=\text{Student}$, then $p.team=\operatorname{getTeamOf}(user,course)$ and $p.supervisor=supervisor$. Otherwise $p.team=\text{None}$ and $p.supervisor=user$ (the argument $supervisor$ is ignored).
    \[
    \mathit{projects}'=\mathit{projects}\cup\{p.id\mapsto p\}.
    \]
    A draft may be incomplete. Persist and audit the change.
    \item output: $out=p$.
    \item exception:
    \begin{itemize}
        \item Unauthorized if $\operatorname{hasAccess}(user,course,\text{Create},course)=\text{false}$.
        \item FeatureDisabled if ProjectProposals is not enabled in the course.
        \item NoTemplate if the course has no template of kind $k$.
        \item NotInTeam if the user is a Student who is not in a team in the course.
        \item InvalidCapacity if $capacity<1$, or the user is a Student and $capacity\ne1$.
        \item InvalidProposal if a field name in $values$ is not in the template.
    \end{itemize}
\end{itemize}

\paragraph{\texttt{updateProposal(user, project, values)}}
\begin{itemize}
    \item transition: Let $c=project.course$. If $\operatorname{hasAccess}(user,c,\text{Update},project)=\text{true}$, $project.status\in\{\text{Draft},\text{ProposalChangesRequested}\}$, and $\operatorname{dom}(values)\subseteq\operatorname{dom}(\mathit{templates}[c][\operatorname{templateKind}(\operatorname{getRole}(project.proposer,c))])$:
    \[
    \mathit{projects}'[project.id].fields=values.
    \]
    Persist and audit the change.
    \item output: None.
    \item exception:
    \begin{itemize}
        \item Unauthorized if $\operatorname{hasAccess}(user,c,\text{Update},project)=\text{false}$.
        \item InvalidStatus if $project.status\notin\{\text{Draft},\text{ProposalChangesRequested}\}$.
        \item InvalidProposal if a field name in $values$ is not in the template.
    \end{itemize}
\end{itemize}

\paragraph{\texttt{submitProposal(user, project)}}
\begin{itemize}
    \item transition: Let $c=project.course$. If $\operatorname{hasAccess}(user,c,\text{Submit},project)=\text{true}$, $\operatorname{isFeatureEnabled}(user,c,\text{ProjectProposals})=\text{true}$, $project.status\in\{\text{Draft},\text{ProposalChangesRequested}\}$, and $\operatorname{missingRequired}(project)=\emptyset$:
    \[
    \mathit{projects}'[project.id].status=\text{ProposalSubmitted}.
    \]
    Call $\operatorname{recordPendingWork}(c,\text{ProjectProposals},project.id,\text{true})$ and $\operatorname{notify}(\text{ProposalSubmitted},project)$. Persist and audit the changes.
    \item output: None.
    \item exception:
    \begin{itemize}
        \item Unauthorized if $\operatorname{hasAccess}(user,c,\text{Submit},project)=\text{false}$.
        \item FeatureDisabled if ProjectProposals is not enabled in the course.
        \item InvalidStatus if $project.status\notin\{\text{Draft},\text{ProposalChangesRequested}\}$.
        \item InvalidProposal if a required field is missing or empty. The exception lists the missing field names.
    \end{itemize}
\end{itemize}

\paragraph{\texttt{withdrawProposal(user, project)}}
\begin{itemize}
    \item transition: Let $c=project.course$. If $\operatorname{hasAccess}(user,c,\text{Withdraw},project)=\text{true}$ and either $project.status\in\{\text{Draft},\text{ProposalSubmitted},\text{ProposalChangesRequested}\}$ or ($project.status=\text{Open}\land project.reserved=\emptyset$):
    \[
    \mathit{projects}'[project.id].status=\text{ProposalWithdrawn}.
    \]
    If the status was ProposalSubmitted, call $\operatorname{recordPendingWork}(c,\text{ProjectProposals},project.id,\text{false})$. Every application to the project is ended (see $\operatorname{endApplication}$). Persist and audit the changes.
    \item output: None.
    \item exception:
    \begin{itemize}
        \item Unauthorized if $\operatorname{hasAccess}(user,c,\text{Withdraw},project)=\text{false}$.
        \item InvalidStatus if the project is in any other status, or has reserved teams.
    \end{itemize}
\end{itemize}

\paragraph{\texttt{deleteProject(user, project)}}
\begin{itemize}
    \item transition: If $\operatorname{hasAccess}(user,project.course,\text{Delete},project)=\text{true}$ and $project.reserved=\emptyset$, remove the project and its applications:
    \[
    \mathit{projects}'=\mathit{projects}\setminus\{project.id\mapsto project\}.
    \]
    If the status was ProposalSubmitted, call $\operatorname{recordPendingWork}(project.course,\text{ProjectProposals},project.id,\text{false})$. Persist and audit the changes.
    \item output: None.
    \item exception:
    \begin{itemize}
        \item Unauthorized if $\operatorname{hasAccess}(user,project.course,\text{Delete},project)=\text{false}$.
        \item HasReservedTeams if $project.reserved\ne\emptyset$.
    \end{itemize}
\end{itemize}

\paragraph{\texttt{reviewProposal(user, project, decision, comment)}}
\begin{itemize}
    \item transition: Let $c=project.course$. If $project.status=\text{ProposalSubmitted}$, $decision\in\{\text{Approve},\text{Reject},\text{RequestChanges}\}$, and $\operatorname{hasAccess}(user,c,decision,project)=\text{true}$:
    \[
    \mathit{projects}'[project.id].status=
    \begin{cases}
        \text{Open}, & decision=\text{Approve},\\
        \text{ProposalRejected}, & decision=\text{Reject},\\
        \text{ProposalChangesRequested}, & decision=\text{RequestChanges},
    \end{cases}
    \]
    and $\mathit{projects}'[project.id].comment=comment$. Call $\operatorname{recordPendingWork}(c,\text{ProjectProposals},project.id,\text{false})$ and notify the proposer of the decision with $\operatorname{notify}(\text{ProposalDecision},project)$. If $decision=\text{Approve}$, also call $\operatorname{notify}(\text{ProjectListed},project)$, and if $\operatorname{isExternalSupervisor}(project.proposer)$ call $\operatorname{verifyExternalSupervisor}(project.proposer,project)$. Persist and audit the changes.
    \item output: None.
    \item exception:
    \begin{itemize}
        \item Unauthorized if $\operatorname{hasAccess}(user,c,decision,project)=\text{false}$.
        \item InvalidStatus if $project.status\ne\text{ProposalSubmitted}$.
        \item InvalidDecision if $decision$ is not Approve, Reject, or RequestChanges.
    \end{itemize}
\end{itemize}

\paragraph{\texttt{setCapacity(user, project, capacity)}}
\begin{itemize}
    \item transition: If $\operatorname{hasAccess}(user,project.course,\text{Update},project)=\text{true}$, $project.status\in\{\text{Draft},\text{ProposalSubmitted},\text{ProposalChangesRequested},\text{Open},\text{Full}\}$, and $capacity\ge\max(1,|project.reserved|)$:
    \[
    \mathit{projects}'[project.id].capacity=capacity.
    \]
    If the status was Open and $capacity=|project.reserved|$ it becomes Full. If it was Full and $capacity>|project.reserved|$ it becomes Open. A status change calls $\operatorname{notify}(\text{ProjectStatusChanged},project)$. Persist and audit the change.
    \item output: None.
    \item exception:
    \begin{itemize}
        \item Unauthorized if $\operatorname{hasAccess}(user,project.course,\text{Update},project)=\text{false}$.
        \item InvalidStatus if the project is in any other status.
        \item InvalidCapacity if $capacity<\max(1,|project.reserved|)$.
    \end{itemize}
\end{itemize}

\paragraph{\texttt{closeSelection(user, course)}}
\begin{itemize}
    \item transition: If $\operatorname{hasAccess}(user,course,\text{Configure},course)=\text{true}$: every project $p$ of the course with $p.status\in\{\text{Open},\text{Full}\}$ becomes Closed, and every application to it in status Applied or SupervisorAccepts is ended (see $\operatorname{endApplication}$). Applications already Agreed or later are unchanged. Persist and audit the changes.
    \item output: None.
    \item exception: Unauthorized if $\operatorname{hasAccess}(user,course,\text{Configure},course)=\text{false}$.
\end{itemize}

\paragraph{\texttt{apply(user, project)}}
\begin{itemize}
    \item transition: Let $c=project.course$ and $t=\operatorname{getTeamOf}(user,c)$. If $\operatorname{hasAccess}(user,c,\text{Apply},project)=\text{true}$, $project.status=\text{Open}$, $t\ne\text{None}$, $t$ has no application to $project$ whose status is not DeclinedOrWithdrawn, and $t$ has no application with status Agreed or later, create an application $a$ with $a.project=project$, $a.team=t$, $a.status=\text{Applied}$:
    \[
    \mathit{applications}'=\mathit{applications}\cup\{a.id\mapsto a\}.
    \]
    An application uses no slot, so the project's status is unchanged. Call $\operatorname{notify}(\text{ApplicationSubmitted},a)$. Persist and audit the changes.
    \item output: $out=a$.
    \item exception:
    \begin{itemize}
        \item Unauthorized if $\operatorname{hasAccess}(user,c,\text{Apply},project)=\text{false}$.
        \item InvalidStatus if $project.status\ne\text{Open}$.
        \item NotInTeam if the user is not in a team in the course.
        \item DuplicateApplication if the team already has an active application to the project.
        \item AlreadyAgreed if the team already has an application with status Agreed or later.
    \end{itemize}
\end{itemize}

\paragraph{\texttt{decideApplication(user, application, decision, comment)}}
\begin{itemize}
    \item transition: Let $p=\mathit{projects}[application.project]$ and $c=p.course$. If $\operatorname{hasAccess}(user,c,decision,application)=\text{true}$ and $decision\in\{\text{Accept},\text{Decline}\}$:
    \begin{itemize}
        \item $decision=\text{Accept}$, allowed when $application.status=\text{Applied}$ and $p.status=\text{Open}$: $\mathit{applications}'[application.id].status=\text{SupervisorAccepts}$ and $\operatorname{reserveSlot}(p,application.team)$.
        \item $decision=\text{Decline}$, allowed when $application.status\in\{\text{Applied},\text{SupervisorAccepts},\text{Agreed},\text{TeamProjectChangesRequested}\}$: $\operatorname{endApplication}(application)$.
    \end{itemize}
    Call $\operatorname{notify}(\text{ApplicationDecision},application)$. The $comment$ (an optional reason) is stored with the notification. Persist and audit the changes.
    \item output: None.
    \item exception:
    \begin{itemize}
        \item Unauthorized if $\operatorname{hasAccess}(user,c,decision,application)=\text{false}$.
        \item InvalidDecision if $decision$ is not Accept or Decline.
        \item InvalidStatus if the application's status does not allow the decision.
        \item NoSlotAvailable if $decision=\text{Accept}$ and $p.status\ne\text{Open}$.
    \end{itemize}
\end{itemize}

\paragraph{\texttt{confirmAgreement(user, application)}}
\begin{itemize}
    \item transition: Let $c$ be the course of the application's project. If $\operatorname{hasAccess}(user,c,\text{Accept},application)=\text{true}$ and $application.status=\text{SupervisorAccepts}$:
    \[
    \mathit{applications}'[application.id].status=\text{Agreed}.
    \]
    The team's slot was already reserved when the supervisor accepted, so no slot changes. Every other application of the same team with status Applied or SupervisorAccepts is ended (see $\operatorname{endApplication}$). Call $\operatorname{notify}(\text{AgreementReached},application)$. All changes are saved in one transaction and audited.
    \item output: None.
    \item exception:
    \begin{itemize}
        \item Unauthorized if $\operatorname{hasAccess}(user,c,\text{Accept},application)=\text{false}$.
        \item InvalidStatus if $application.status\ne\text{SupervisorAccepts}$ (the team cannot accept before the supervisor).
    \end{itemize}
\end{itemize}

\paragraph{\texttt{withdrawApplication(user, application)}}
\begin{itemize}
    \item transition: Let $c$ be the course of the application's project. If $\operatorname{hasAccess}(user,c,\text{Withdraw},application)=\text{true}$ and $application.status\in\{\text{Applied},\text{SupervisorAccepts},\text{Agreed},\text{TeamProjectChangesRequested},\text{Approved}\}$: $\operatorname{endApplication}(application)$. Call $\operatorname{notify}(\text{ApplicationWithdrawn},application)$. Persist and audit the changes. Withdrawing an Approved application is how a team that dissolves after approval releases its slot.
    \item output: None.
    \item exception:
    \begin{itemize}
        \item Unauthorized if $\operatorname{hasAccess}(user,c,\text{Withdraw},application)=\text{false}$.
        \item InvalidStatus if the application is in any other status, including while a review of it is in progress.
    \end{itemize}
\end{itemize}

\paragraph{\texttt{updateApplicationStatus(application, status)}}
\begin{itemize}
    \item transition: Called only by M\_ApprovalWorkflow. If $\operatorname{validApplicationTransition}(application.status,status)$:
    \[
    \mathit{applications}'[application.id].status=status.
    \]
    If $status=\text{TeamProjectRejected}$ and $\operatorname{holdsSlot}(application.status)$, call $\operatorname{releaseSlot}(\mathit{projects}[application.project],application.team)$. The changes are saved in the caller's transaction and audited.
    \item output: None.
    \item exception: InvalidTransition if $\neg\operatorname{validApplicationTransition}(application.status,status)$.
\end{itemize}

\paragraph{\texttt{getProject(user, project)}}
\begin{itemize}
    \item transition: None.
    \item output: If $\operatorname{hasAccess}(user,project.course,\text{Read},project)=\text{true}$:
    \[
    out=\mathit{projects}[project.id].
    \]
    \item exception: Unauthorized if $\operatorname{hasAccess}(user,project.course,\text{Read},project)=\text{false}$.
\end{itemize}

\paragraph{\texttt{getApplication(user, application)}}
\begin{itemize}
    \item transition: None.
    \item output: If $\operatorname{hasAccess}(user,c,\text{Read},application)=\text{true}$, where $c$ is the course of the application's project:
    \[
    out=\mathit{applications}[application.id].
    \]
    \item exception: Unauthorized if $\operatorname{hasAccess}(user,c,\text{Read},application)=\text{false}$.
\end{itemize}

\paragraph{\texttt{listProjects(user, course, text, statuses, sortBy)}}
\begin{itemize}
    \item transition: None.
    \item output: If $\operatorname{hasAccess}(user,course,\text{Read},course)=\text{true}$ and $sortBy\in\{\text{title},\text{remaining}\}$, the list
    \[
    [p\mid p\in\mathit{projects},\ p.course=course\land\operatorname{hasAccess}(user,course,\text{Read},p)\land(statuses=\emptyset\lor p.status\in statuses)\land\operatorname{matches}(p,text)],
    \]
    ordered by title ascending if $sortBy=\text{title}$, or by $\operatorname{remaining}(p)$ descending if $sortBy=\text{remaining}$. Ties are ordered by identifier.
    \item exception:
    \begin{itemize}
        \item Unauthorized if $\operatorname{hasAccess}(user,course,\text{Read},course)=\text{false}$.
        \item InvalidQuery if $sortBy$ is not title or remaining.
    \end{itemize}
\end{itemize}

\paragraph{\texttt{getApplications(user, project)}}
\begin{itemize}
    \item transition: None.
    \item output: If $\operatorname{hasAccess}(user,project.course,\text{Read},project)=\text{true}$, the applications to the project that the user may read:
    \[
    out=[a\mid a\in\mathit{applications},\ a.project=project\land\operatorname{hasAccess}(user,project.course,\text{Read},a)].
    \]
    \item exception: Unauthorized if $\operatorname{hasAccess}(user,project.course,\text{Read},project)=\text{false}$.
\end{itemize}

\paragraph{\texttt{getTeamApplications(user, team)}}
\begin{itemize}
    \item transition: None.
    \item output: If $\operatorname{hasAccess}(user,team.course,\text{Read},team)=\text{true}$, the applications of the team that the user may read:
    \[
    out=[a\mid a\in\mathit{applications},\ a.team=team\land\operatorname{hasAccess}(user,team.course,\text{Read},a)].
    \]
    \item exception: Unauthorized if $\operatorname{hasAccess}(user,team.course,\text{Read},team)=\text{false}$.
\end{itemize}

\subsubsection{Local Functions}

\begin{description}
    \item[templateKind(role: UserRole): String] Returns Student for the Student role and Supervisor for every other role.
    \item[missingRequired(project: Project): $\operatorname{Set}\langle\text{String}\rangle$] Returns the names of fields marked required in the template for the proposer's kind that are absent or empty in $project.fields$.
    \item[remaining(project: Project): $\mathbb{N}$] Returns $project.capacity-|project.reserved|$.
    \item[matches(project: Project, text: String): bool] Returns true if $text=\epsilon$ or $text$ occurs, ignoring case, in some value of $project.fields$.
    \item[holdsSlot(status: TeamProjectStatus): bool] Returns true for SupervisorAccepts, Agreed, TeamProjectSubmitted, TeamProjectUnderReview, TeamProjectChangesRequested, Approved, and ApprovedChangeUnderReview.
    \item[reserveSlot(project: Project, team: Team): --] Adds $team.id$ to $project.reserved$. If $|project.reserved|$ now equals $project.capacity$ the project becomes Full, and a status change calls $\operatorname{notify}(\text{ProjectStatusChanged},project)$.
    \item[releaseSlot(project: Project, team: Team): --] Removes $team.id$ from $project.reserved$. If the project was Full it becomes Open and $\operatorname{notify}(\text{ProjectStatusChanged},project)$ is called.
    \item[endApplication(application: ProjectApplication): --] Sets the application's status to DeclinedOrWithdrawn. If $\operatorname{holdsSlot}$ was true for its previous status, calls $\operatorname{releaseSlot}$ for its project and team.
    \item[validApplicationTransition(from, to: TeamProjectStatus): bool] Returns true exactly for these changes: TeamProjectSubmitted to TeamProjectUnderReview, TeamProjectChangesRequested, TeamProjectRejected, or Approved; TeamProjectUnderReview to TeamProjectChangesRequested, TeamProjectRejected, or Approved; Agreed or TeamProjectChangesRequested to TeamProjectSubmitted; Approved to ApprovedChangeUnderReview; ApprovedChangeUnderReview to Approved. (Changes into Applied, SupervisorAccepts, Agreed, and DeclinedOrWithdrawn are made only by this module's own routines.)
    \item[isExternalSupervisor(user: User): bool] Returns true if the user's linked sign-in account is with the ExternalProvider.
\end{description}

\newpage

\section[MIS of Students and Teams Module]{MIS of Students and Teams Module (\moduleid{M_StudentsTeams})}
\label{M_StudentsTeams}

\subsection{Module}

This module manages how students form teams. It stores and shows student and team profiles, creates teams, handles invitations and join requests, records member roles and the assigned TA, and finds students and teams that are looking for each other. It checks a team against the course's team rules, which are read from the course settings, and reports any rule the team breaks.

\subsection{Uses}

\begin{itemize}
    \item M\_Authorization
    \item M\_CourseManagement
    \item M\_Notifications
    \item M\_Persistence
    \item M\_AuditLog
\end{itemize}

\subsection{Syntax}

\subsubsection{Exported Constants}
None.

\subsubsection{Exported Access Programs}

\begingroup
\small
\setlength{\tabcolsep}{4pt}
\renewcommand{\arraystretch}{1.2}
\begin{longtable}{>{\raggedright\arraybackslash}p{0.24\textwidth}
                  >{\raggedright\arraybackslash}p{0.34\textwidth}
                  >{\raggedright\arraybackslash}p{0.16\textwidth}
                  >{\raggedright\arraybackslash}p{0.16\textwidth}}
\toprule
\textbf{Name} & \textbf{In} & \textbf{Out} & \textbf{Exceptions} \\
\midrule
\endfirsthead
\toprule
\textbf{Name} & \textbf{In} & \textbf{Out} & \textbf{Exceptions} \\
\midrule
\endhead
\endfoot
\texttt{saveStudentProfile} & user: User, course: Course, profile: StudentProfile & -- & Unauthorized, InvalidProfile \\
\texttt{getStudentProfile} & user: User, course: Course, student: User & StudentProfile & Unauthorized, ProfileNotFound \\
\texttt{deleteStudentProfile} & user: User, course: Course & -- & Unauthorized, ProfileNotFound, InTeam \\
\texttt{createTeam} & user: User, course: Course, profile: TeamProfile & Team & Unauthorized, FeatureDisabled, NoProfile, AlreadyInTeam, InvalidProfile \\
\texttt{saveTeamProfile} & user: User, team: Team, profile: TeamProfile & -- & Unauthorized, InvalidProfile \\
\texttt{setMembersNeeded} & user: User, team: Team, count: $\mathbb{N}$ & -- & Unauthorized \\
\texttt{getTeam} & user: User, team: Team & Team & Unauthorized \\
\texttt{getTeamOf} & user: User, course: Course & $\text{Team}\cup\{\text{None}\}$ & Unauthorized \\
\texttt{listTeams} & user: User, course: Course, recruitingOnly: bool & $\operatorname{List}\langle\text{Team}\rangle$ & Unauthorized \\
\texttt{listSeeking} & user: User, course: Course & $\operatorname{List}\langle\text{User}\rangle$ & Unauthorized \\
\texttt{listUnteamed} & user: User, course: Course, enrolled: $\operatorname{List}\langle\text{User}\rangle$ & $\operatorname{List}\langle\text{User}\rangle$ & Unauthorized \\
\texttt{listNoncompliantTeams} & user: User, course: Course & $\operatorname{List}\langle\text{Team}\rangle$ & Unauthorized \\
\texttt{invite} & user: User, team: Team, student: User & Offer & Unauthorized, FeatureDisabled, NoProfile, AlreadyInTeam, TeamFull, DuplicateOffer \\
\texttt{requestToJoin} & user: User, team: Team & Offer & Unauthorized, FeatureDisabled, NoProfile, AlreadyInTeam, TeamFull, DuplicateOffer \\
\texttt{respondToOffer} & user: User, offer: Offer, accept: bool & -- & Unauthorized, OfferNotFound, AlreadyInTeam, TeamFull \\
\texttt{leaveTeam} & user: User, team: Team & -- & Unauthorized, NotAMember \\
\texttt{removeMember} & user: User, team: Team, member: User & -- & Unauthorized, NotAMember, CannotRemovePointOfContact \\
\texttt{setPointOfContact} & user: User, team: Team, member: User & -- & Unauthorized, NotAMember \\
\texttt{dissolveTeam} & user: User, team: Team & -- & Unauthorized \\
\texttt{assignTA} & user: User, team: Team, ta: User & -- & Unauthorized, InvalidTA \\
\texttt{validateTeam} & user: User, team: Team & $\operatorname{Set}\langle\text{RuleViolation}\rangle$ & Unauthorized \\
\bottomrule
\caption{Exported access programs of the Students and Teams module}
\label{TblStudentsTeamsAccessPrograms}
\end{longtable}
\endgroup

\subsection{Semantics}

\subsubsection{State Variables}

\begin{itemize}
    \item $\mathit{profiles}:\operatorname{Map}\langle\text{Course},\operatorname{Map}\langle\text{User},\text{StudentProfile}\rangle\rangle$ - student profiles of each course, initially loaded from M\_Persistence.

    \item $\mathit{teams}:\operatorname{Map}\langle\text{String},\text{Team}\rangle$ - teams indexed by identifier, initially loaded from M\_Persistence.

    \item $\mathit{offers}:\operatorname{Map}\langle\text{String},\text{Offer}\rangle$ - pending invitations and join requests indexed by identifier, initially loaded from M\_Persistence.
\end{itemize}

\subsubsection{Environment Variables}

Persistent profile, team, and offer records.

\subsubsection{Assumptions}

\begin{itemize}
    \item Users are authenticated through M\_Authentication, and are enrolled in the course, before any routine is called.
    \item M\_Authorization enforces the SRS visibility rules through \texttt{hasAccess}.
    \item The caller has re-verified the user, as required by \texttt{FR\_ReverifyUsers}, before calling \texttt{deleteStudentProfile}, \texttt{leaveTeam}, and \texttt{dissolveTeam}.
    \item A user is in at most one team per course.
    \item A team's rules are the course settings \texttt{minTeamSize}, \texttt{maxTeamSize}, and \texttt{allowedCombinations}, set through M\_CourseManagement. A rule that is not set is not enforced.
    \item The caller has ended every application of a team that holds a project slot (with \texttt{withdrawApplication} in M\_Projects) before the team's last member leaves or the team is dissolved, because this module cannot see applications.
    \item Pending offers count as pending work of the TeamFormation feature, so M\_CourseManagement refuses to turn the feature off while any exist.
    \item $\operatorname{notify}(event,record)$ is the M\_Notifications routine that decides the recipients, stores the notification in the caller's transaction, and delivers it later. M\_Persistence keeps the \texttt{updated\_at} value of each record.
\end{itemize}

\subsubsection{Access Routine Semantics}

\paragraph{\texttt{saveStudentProfile(user, course, profile)}}
\begin{itemize}
    \item transition: Let $a=\text{Update}$ if $user\in\operatorname{dom}(\mathit{profiles}[course])$, otherwise $a=\text{Create}$. If $\operatorname{hasAccess}(user,course,a,profile)=\text{true}$ and $\operatorname{validProfile}(profile)$:
    \[
    \mathit{profiles}'[course]=\mathit{profiles}[course]\oplus\{user\mapsto profile\}.
    \]
    Persist and audit the change.
    \item output: None.
    \item exception:
    \begin{itemize}
        \item Unauthorized if $\operatorname{hasAccess}(user,course,a,profile)=\text{false}$.
        \item InvalidProfile if $\neg\operatorname{validProfile}(profile)$.
    \end{itemize}
\end{itemize}

\paragraph{\texttt{getStudentProfile(user, course, student)}}
\begin{itemize}
    \item transition: None.
    \item output: If $student\in\operatorname{dom}(\mathit{profiles}[course])$ and $\operatorname{hasAccess}(user,course,\text{Read},\mathit{profiles}[course][student])=\text{true}$:
    \[
    out=\mathit{profiles}[course][student].
    \]
    \item exception:
    \begin{itemize}
        \item ProfileNotFound if $student\notin\operatorname{dom}(\mathit{profiles}[course])$.
        \item Unauthorized if $\operatorname{hasAccess}(user,course,\text{Read},\mathit{profiles}[course][student])=\text{false}$.
    \end{itemize}
\end{itemize}

\paragraph{\texttt{deleteStudentProfile(user, course)}}
\begin{itemize}
    \item transition: If $user\in\operatorname{dom}(\mathit{profiles}[course])$, $\operatorname{teamOf}(user,course)=\text{None}$, and $\operatorname{hasAccess}(user,course,\text{Delete},\mathit{profiles}[course][user])=\text{true}$:
    \[
    \mathit{profiles}'[course]=\mathit{profiles}[course]\setminus\{user\mapsto\mathit{profiles}[course][user]\}.
    \]
    Pending offers involving the user are removed (see $\operatorname{dropOffers}$). Persist and audit the changes.
    \item output: None.
    \item exception:
    \begin{itemize}
        \item ProfileNotFound if $user\notin\operatorname{dom}(\mathit{profiles}[course])$.
        \item InTeam if the user is in a team in the course.
        \item Unauthorized if $\operatorname{hasAccess}(user,course,\text{Delete},\mathit{profiles}[course][user])=\text{false}$.
    \end{itemize}
\end{itemize}

\paragraph{\texttt{createTeam(user, course, profile)}}
\begin{itemize}
    \item transition: If $\operatorname{hasAccess}(user,course,\text{Create},course)=\text{true}$, $\operatorname{isFeatureEnabled}(user,course,\text{TeamFormation})=\text{true}$, $user\in\operatorname{dom}(\mathit{profiles}[course])$, $\operatorname{teamOf}(user,course)=\text{None}$, and $\operatorname{validTeamProfile}(profile)$, create a team $t$ with $t.course=course$, $t.members=\{user\mapsto\text{PointOfContact}\}$, $t.membersNeeded=0$, $t.ta=\text{None}$, and $t.profile=profile$:
    \[
    \mathit{teams}'=\mathit{teams}\cup\{t.id\mapsto t\}.
    \]
    Pending offers involving $user$ in the course are removed (see $\operatorname{dropOffers}$). Persist and audit the changes.
    \item output: $out=t$.
    \item exception:
    \begin{itemize}
        \item Unauthorized if $\operatorname{hasAccess}(user,course,\text{Create},course)=\text{false}$.
        \item FeatureDisabled if TeamFormation is not enabled in the course.
        \item NoProfile if the user has no student profile in the course.
        \item AlreadyInTeam if the user is already in a team in the course.
        \item InvalidProfile if $\neg\operatorname{validTeamProfile}(profile)$.
    \end{itemize}
\end{itemize}

\paragraph{\texttt{saveTeamProfile(user, team, profile)}}
\begin{itemize}
    \item transition: If $\operatorname{hasAccess}(user,team.course,\text{Update},team)=\text{true}$ and $\operatorname{validTeamProfile}(profile)$:
    \[
    \mathit{teams}'[team.id].profile=profile.
    \]
    Persist and audit the change.
    \item output: None.
    \item exception:
    \begin{itemize}
        \item Unauthorized if $\operatorname{hasAccess}(user,team.course,\text{Update},team)=\text{false}$.
        \item InvalidProfile if $\neg\operatorname{validTeamProfile}(profile)$.
    \end{itemize}
\end{itemize}

\paragraph{\texttt{setMembersNeeded(user, team, count)}}
\begin{itemize}
    \item transition: If $\operatorname{hasAccess}(user,team.course,\text{Update},team)=\text{true}$:
    \[
    \mathit{teams}'[team.id].membersNeeded=count.
    \]
    A team with $count>0$ is recruiting. Persist and audit the change.
    \item output: None.
    \item exception: Unauthorized if $\operatorname{hasAccess}(user,team.course,\text{Update},team)=\text{false}$.
\end{itemize}

\paragraph{\texttt{getTeam(user, team)}}
\begin{itemize}
    \item transition: None.
    \item output: If $\operatorname{hasAccess}(user,team.course,\text{Read},team)=\text{true}$:
    \[
    out=\mathit{teams}[team.id].
    \]
    \item exception: Unauthorized if $\operatorname{hasAccess}(user,team.course,\text{Read},team)=\text{false}$.
\end{itemize}

\paragraph{\texttt{getTeamOf(user, course)}}
\begin{itemize}
    \item transition: None.
    \item output: If $\operatorname{hasAccess}(user,course,\text{Read},course)=\text{true}$:
    \[
    out=\operatorname{teamOf}(user,course).
    \]
    \item exception: Unauthorized if $\operatorname{hasAccess}(user,course,\text{Read},course)=\text{false}$.
\end{itemize}

\paragraph{\texttt{listTeams(user, course, recruitingOnly)}}
\begin{itemize}
    \item transition: None.
    \item output: If $\operatorname{hasAccess}(user,course,\text{Read},course)=\text{true}$, the teams of the course that the user may read, ordered by identifier:
    \[
    out=[t\mid t\in\mathit{teams},\ t.course=course\land\operatorname{hasAccess}(user,course,\text{Read},t)\land(\neg recruitingOnly\lor t.membersNeeded>0)].
    \]
    \item exception: Unauthorized if $\operatorname{hasAccess}(user,course,\text{Read},course)=\text{false}$.
\end{itemize}

\paragraph{\texttt{listSeeking(user, course)}}
\begin{itemize}
    \item transition: None.
    \item output: If $\operatorname{hasAccess}(user,course,\text{Read},course)=\text{true}$, the students who are looking for a team, ordered by identifier:
    \[
    out=[u\mid u\in\operatorname{dom}(\mathit{profiles}[course]),\ \mathit{profiles}[course][u].seekingTeam\land\operatorname{teamOf}(u,course)=\text{None}\land\operatorname{hasAccess}(user,course,\text{Read},\mathit{profiles}[course][u])].
    \]
    \item exception: Unauthorized if $\operatorname{hasAccess}(user,course,\text{Read},course)=\text{false}$.
\end{itemize}

\paragraph{\texttt{listUnteamed(user, course, enrolled)}}
\begin{itemize}
    \item transition: None.
    \item output: If $\operatorname{hasAccess}(user,course,\text{ViewReport},course)=\text{true}$, the enrolled students who are on no team, in the order of $enrolled$:
    \[
    out=[u\mid u\in enrolled,\ \operatorname{teamOf}(u,course)=\text{None}].
    \]
    The caller supplies $enrolled$ from M\_Enrollment, which this module does not use.
    \item exception: Unauthorized if $\operatorname{hasAccess}(user,course,\text{ViewReport},course)=\text{false}$.
\end{itemize}

\paragraph{\texttt{listNoncompliantTeams(user, course)}}
\begin{itemize}
    \item transition: None.
    \item output: If $\operatorname{hasAccess}(user,course,\text{ViewReport},course)=\text{true}$, the teams of the course that break at least one rule, ordered by identifier, so that staff can review them after the rules change. The teams are not changed or invalidated:
    \[
    out=[t\mid t\in\mathit{teams},\ t.course=course\land\operatorname{violations}(user,t)\ne\emptyset].
    \]
    \item exception: Unauthorized if $\operatorname{hasAccess}(user,course,\text{ViewReport},course)=\text{false}$.
\end{itemize}

\paragraph{\texttt{invite(user, team, student)}}
\begin{itemize}
    \item transition: Let $c=team.course$ and let $o$ be a new offer with $o.team=team$, $o.user=student$, and $o.kind=\text{Invitation}$. If $\operatorname{hasAccess}(user,c,\text{AddMember},team)=\text{true}$, $\operatorname{isFeatureEnabled}(user,c,\text{TeamFormation})=\text{true}$, $student\in\operatorname{dom}(\mathit{profiles}[c])$, $\operatorname{teamOf}(student,c)=\text{None}$, $\neg\operatorname{isFull}(user,team)$, and no offer exists for the same team and student:
    \[
    \mathit{offers}'=\mathit{offers}\cup\{o.id\mapsto o\}.
    \]
    Call $\operatorname{recordPendingWork}(c,\text{TeamFormation},o.id,\text{true})$ and $\operatorname{notify}(\text{TeamInvitation},o)$. Persist and audit the changes.
    \item output: $out=o$.
    \item exception:
    \begin{itemize}
        \item Unauthorized if $\operatorname{hasAccess}(user,c,\text{AddMember},team)=\text{false}$.
        \item FeatureDisabled if TeamFormation is not enabled in the course.
        \item NoProfile if $student$ has no student profile in the course.
        \item AlreadyInTeam if $student$ is already in a team in the course.
        \item TeamFull if $\operatorname{isFull}(user,team)$.
        \item DuplicateOffer if an offer already exists for the team and student.
    \end{itemize}
\end{itemize}

\paragraph{\texttt{requestToJoin(user, team)}}
\begin{itemize}
    \item transition: Let $c=team.course$ and let $o$ be a new offer with $o.team=team$, $o.user=user$, and $o.kind=\text{Request}$. If $\operatorname{hasAccess}(user,c,\text{Create},o)=\text{true}$ and the other conditions of \texttt{invite} hold with $user$ in place of $student$:
    \[
    \mathit{offers}'=\mathit{offers}\cup\{o.id\mapsto o\}.
    \]
    Call $\operatorname{recordPendingWork}(c,\text{TeamFormation},o.id,\text{true})$ and $\operatorname{notify}(\text{JoinRequest},o)$. Persist and audit the changes.
    \item output: $out=o$.
    \item exception: The exceptions of \texttt{invite}, with Unauthorized if $\operatorname{hasAccess}(user,c,\text{Create},o)=\text{false}$.
\end{itemize}

\paragraph{\texttt{respondToOffer(user, offer, accept)}}
\begin{itemize}
    \item transition: Let $t=\mathit{teams}[offer.team]$ and $c=t.course$. If $offer.id\in\operatorname{dom}(\mathit{offers})$, $\operatorname{hasAccess}(user,c,\operatorname{responseAction}(offer,accept),offer)=\text{true}$, and, when $accept$, $\operatorname{teamOf}(offer.user,c)=\text{None}\land\neg\operatorname{isFull}(user,t)$:
    \[
    \mathit{offers}'=\mathit{offers}\setminus\{offer.id\mapsto offer\}.
    \]
    Call $\operatorname{recordPendingWork}(c,\text{TeamFormation},offer.id,\text{false})$ and $\operatorname{notify}(\text{OfferDecision},offer)$. If $accept$:
    \[
    \mathit{teams}'[t.id].members=t.members\oplus\{offer.user\mapsto\text{Member}\},
    \]
    and every other pending offer involving $offer.user$ in the course is removed (see $\operatorname{dropOffers}$). Persist and audit the changes.
    \item output: None.
    \item exception:
    \begin{itemize}
        \item OfferNotFound if $offer.id\notin\operatorname{dom}(\mathit{offers})$.
        \item Unauthorized if $\operatorname{hasAccess}(user,c,\operatorname{responseAction}(offer,accept),offer)=\text{false}$.
        \item AlreadyInTeam if $accept$ and $offer.user$ is already in a team in the course.
        \item TeamFull if $accept$ and $\operatorname{isFull}(user,t)$.
    \end{itemize}
\end{itemize}

\paragraph{\texttt{leaveTeam(user, team)}}
\begin{itemize}
    \item transition: If $\operatorname{hasAccess}(user,team.course,\text{LeaveTeam},team)=\text{true}$ and $user\in\operatorname{dom}(team.members)$, remove the user:
    \[
    \mathit{teams}'[team.id].members=team.members\setminus\{user\mapsto team.members[user]\}.
    \]
    If the user was the PointOfContact and members remain, $\operatorname{successor}(team)$ becomes PointOfContact. If no members remain, the team and its pending offers are removed (see $\operatorname{dropOffers}$). Call $\operatorname{notify}(\text{MemberLeft},team)$. Persist and audit the changes.
    \item output: None.
    \item exception:
    \begin{itemize}
        \item Unauthorized if $\operatorname{hasAccess}(user,team.course,\text{LeaveTeam},team)=\text{false}$.
        \item NotAMember if $user\notin\operatorname{dom}(team.members)$.
    \end{itemize}
\end{itemize}

\paragraph{\texttt{removeMember(user, team, member)}}
\begin{itemize}
    \item transition: If $\operatorname{hasAccess}(user,team.course,\text{RemoveMember},team)=\text{true}$, $member\in\operatorname{dom}(team.members)$, and $team.members[member]=\text{Member}$:
    \[
    \mathit{teams}'[team.id].members=team.members\setminus\{member\mapsto\text{Member}\}.
    \]
    Call $\operatorname{notify}(\text{MemberRemoved},team)$. Persist and audit the changes.
    \item output: None.
    \item exception:
    \begin{itemize}
        \item Unauthorized if $\operatorname{hasAccess}(user,team.course,\text{RemoveMember},team)=\text{false}$.
        \item NotAMember if $member\notin\operatorname{dom}(team.members)$.
        \item CannotRemovePointOfContact if $team.members[member]=\text{PointOfContact}$.
    \end{itemize}
\end{itemize}

\paragraph{\texttt{setPointOfContact(user, team, member)}}
\begin{itemize}
    \item transition: If $\operatorname{hasAccess}(user,team.course,\text{Update},team)=\text{true}$ and $member\in\operatorname{dom}(team.members)$, $member$ becomes PointOfContact and the previous PointOfContact becomes Member. Persist and audit the change.
    \item output: None.
    \item exception:
    \begin{itemize}
        \item Unauthorized if $\operatorname{hasAccess}(user,team.course,\text{Update},team)=\text{false}$.
        \item NotAMember if $member\notin\operatorname{dom}(team.members)$.
    \end{itemize}
\end{itemize}

\paragraph{\texttt{dissolveTeam(user, team)}}
\begin{itemize}
    \item transition: If $\operatorname{hasAccess}(user,team.course,\text{DissolveTeam},team)=\text{true}$, remove the team from $\mathit{teams}$, so its members have no team, and remove its pending offers (see $\operatorname{dropOffers}$). Call $\operatorname{notify}(\text{TeamDissolved},team)$. Persist and audit the changes.
    \item output: None.
    \item exception: Unauthorized if $\operatorname{hasAccess}(user,team.course,\text{DissolveTeam},team)=\text{false}$.
\end{itemize}

\paragraph{\texttt{assignTA(user, team, ta)}}
\begin{itemize}
    \item transition: If $\operatorname{hasAccess}(user,team.course,\text{Assign},team)=\text{true}$ and $\operatorname{getRole}(ta,team.course)=\text{TA}$:
    \[
    \mathit{teams}'[team.id].ta=ta.
    \]
    Call $\operatorname{notify}(\text{TAAssigned},team)$. Persist and audit the changes.
    \item output: None.
    \item exception:
    \begin{itemize}
        \item Unauthorized if $\operatorname{hasAccess}(user,team.course,\text{Assign},team)=\text{false}$.
        \item InvalidTA if $ta$ does not have the TA role in the course.
    \end{itemize}
\end{itemize}

\paragraph{\texttt{validateTeam(user, team)}}
\begin{itemize}
    \item transition: None.
    \item output: If $\operatorname{hasAccess}(user,team.course,\text{Read},team)=\text{true}$:
    \[
    out=\operatorname{violations}(user,team).
    \]
    An empty set means the team satisfies every rule. M\_ApprovalWorkflow uses this check before a submission is sent.
    \item exception: Unauthorized if $\operatorname{hasAccess}(user,team.course,\text{Read},team)=\text{false}$.
\end{itemize}

\subsubsection{Local Functions}

\begin{description}
    \item[teamOf(user: User, course: Course): $\text{Team}\cup\{\text{None}\}$] Returns the team of the course that has $user$ as a member, or None.
    \item[rules(user: User, course: Course): $(\mathbb{N},\ \mathbb{N}\cup\{\text{None}\},\ \operatorname{Set}\langle\operatorname{Set}\langle\text{Program}\rangle\rangle)$] Returns $(\mathit{min},\mathit{max},\mathit{combos})$ from the course settings \texttt{minTeamSize}, \texttt{maxTeamSize}, and \texttt{allowedCombinations}, read with $\operatorname{getSettings}(user,course)$. Unset values are $1$, None (no maximum), and $\emptyset$ (any combination).
    \item[isFull(user: User, team: Team): bool] Returns true if $\mathit{max}\ne\text{None}\land|team.members|\ge\mathit{max}$, with $\mathit{max}$ from $\operatorname{rules}(user,team.course)$.
    \item[violations(user: User, team: Team): $\operatorname{Set}\langle\text{RuleViolation}\rangle$] With $(\mathit{min},\mathit{max},\mathit{combos})=\operatorname{rules}(user,team.course)$, $m=|team.members|$, and $P$ the set of programs of the members' student profiles, returns the set containing TooFewMembers if $m<\mathit{min}$; TooManyMembers if $\mathit{max}\ne\text{None}\land m>\mathit{max}$; and ProgramCombinationNotAllowed if $\mathit{combos}\ne\emptyset\land P\notin\mathit{combos}$.
    \item[validProfile(profile: StudentProfile): bool] Returns true if the profile has valid values for the fields the SRS requires, including its program.
    \item[validTeamProfile(profile: TeamProfile): bool] Returns true if the profile has valid values for the fields the SRS requires.
    \item[responseAction(offer: Offer, accept: bool): Action] Returns Accept or Decline if $offer.kind=\text{Invitation}$ (the invited student responds). Returns AddMember if $offer.kind=\text{Request}\land accept$, and Decline if $offer.kind=\text{Request}\land\neg accept$ (a team member responds).
    \item[dropOffers(...): --] Removes the given pending offers from $\mathit{offers}$ and calls $\operatorname{recordPendingWork}(course,\text{TeamFormation},id,\text{false})$ for each one removed.
    \item[successor(team: Team): User] Returns the remaining member of the team who joined earliest.
\end{description}

\newpage

\section[MIS of Recommendation Module]{MIS of Recommendation Module (\moduleid{M_Recommendation})}
\label{M_Recommendation}

\subsection{Module}

This module suggests teammates, teams, and projects. It reads the profiles and projects the user may see through the modules that own them, and ranks them by how well they match the user's team profile (or student profile, if the user has no team). The matching method is the secret of the module, so it can be changed, extended to more criteria, or removed without changing any other module. The module keeps no state of its own.

\subsection{Uses}

\begin{itemize}
    \item M\_Authorization
    \item M\_StudentsTeams
    \item M\_Projects
\end{itemize}

\subsection{Syntax}

\subsubsection{Exported Constants}
None.

\subsubsection{Exported Access Programs}

\begingroup
\small
\setlength{\tabcolsep}{4pt}
\renewcommand{\arraystretch}{1.2}
\begin{longtable}{>{\raggedright\arraybackslash}p{0.24\textwidth}
                  >{\raggedright\arraybackslash}p{0.34\textwidth}
                  >{\raggedright\arraybackslash}p{0.16\textwidth}
                  >{\raggedright\arraybackslash}p{0.16\textwidth}}
\toprule
\textbf{Name} & \textbf{In} & \textbf{Out} & \textbf{Exceptions} \\
\midrule
\endfirsthead
\toprule
\textbf{Name} & \textbf{In} & \textbf{Out} & \textbf{Exceptions} \\
\midrule
\endhead
\endfoot
\texttt{recommendTeammates} & user: User, course: Course, k: $\mathbb{N}$ & $\operatorname{List}\langle\text{User}\rangle$ & Unauthorized, NoProfile \\
\texttt{recommendTeams} & user: User, course: Course, k: $\mathbb{N}$ & $\operatorname{List}\langle\text{Team}\rangle$ & Unauthorized, NoProfile \\
\texttt{recommendProjects} & user: User, course: Course, k: $\mathbb{N}$ & $\operatorname{List}\langle\text{Project}\rangle$ & Unauthorized, NoProfile \\
\bottomrule
\caption{Exported access programs of the Recommendation module}
\label{TblRecommendationAccessPrograms}
\end{longtable}
\endgroup

\subsection{Semantics}

\subsubsection{State Variables}
None.

\subsubsection{Environment Variables}
None.

\subsubsection{Assumptions}

\begin{itemize}
    \item Users are authenticated through M\_Authentication before any routine is called.
    \item M\_StudentsTeams and M\_Projects return only the records the user may read, so this module does not check access to individual records.
    \item A student without a team has a student profile before asking for recommendations.
\end{itemize}

\subsubsection{Access Routine Semantics}

In all three routines, the \emph{seeker} is $\operatorname{getTeamOf}(user,course)$ with its team profile if the user is in a team, and otherwise the user's student profile (from $\operatorname{getStudentProfile}$).

\paragraph{\texttt{recommendTeammates(user, course, k)}}
\begin{itemize}
    \item transition: None.
    \item output: If $\operatorname{hasAccess}(user,course,\text{Read},course)=\text{true}$, let $C=\operatorname{listSeeking}(user,course)\setminus\{user\}$ (the students looking for a team). Then $out$ is a list of users such that:
    \[
    \begin{aligned}
    &out\text{ contains only elements of }C,\text{ each at most once},\quad|out|\le k,\\
    &\forall c\in out:\operatorname{match}(\mathit{seeker},c)>0,\\
    &out\text{ is ordered from highest to lowest }\operatorname{match}(\mathit{seeker},c),\text{ ties by identifier},\\
    &|out|=\min\left(k,\ |\{c\in C\mid\operatorname{match}(\mathit{seeker},c)>0\}|\right).
    \end{aligned}
    \]
    \item exception:
    \begin{itemize}
        \item Unauthorized if $\operatorname{hasAccess}(user,course,\text{Read},course)=\text{false}$.
        \item NoProfile if the user has no team and no student profile.
    \end{itemize}
\end{itemize}

\paragraph{\texttt{recommendTeams(user, course, k)}}
\begin{itemize}
    \item transition: None.
    \item output: As \texttt{recommendTeammates}, with $C=\operatorname{listTeams}(user,course,\text{true})$ (the teams that are recruiting), the user's own team removed, and teams matched on their team profiles.
    \item exception: As \texttt{recommendTeammates}.
\end{itemize}

\paragraph{\texttt{recommendProjects(user, course, k)}}
\begin{itemize}
    \item transition: None.
    \item output: As \texttt{recommendTeammates}, with $C=\operatorname{listProjects}(user,course,\epsilon,\{\text{Open}\},\text{title})$ (the Open projects the user may see) and projects matched on their fields.
    \item exception: As \texttt{recommendTeammates}.
\end{itemize}

\subsubsection{Local Functions}

\begin{description}
    \item[match(x: Value, c: Value): $[0,1]$] Returns how well candidate $c$ fits the seeker $x$, where 0 means no match. It is based on shared skills, interests, experience, and project preferences, and depends only on the profile or project information of its arguments, so equal inputs always give equal results. This function is the secret of the module. Changing the matching method changes only this function.
\end{description}

\newpage

\section[MIS of Approval Workflow Module]{MIS of Approval Workflow Module (\moduleid{M_ApprovalWorkflow})}
\label{M_ApprovalWorkflow}

\subsection{Module}

This module runs each course's approval of team projects, in the course's approval mode (set in the course settings). In mode InPlatform, it checks a team's submission before it is sent, holds it for review, and records the decision of an instructor or TA, and it handles changes to approved projects while keeping the existing approval in force until the change is accepted. In mode External, it records the outcome that an instructor or TA enters by hand. In mode Integrated, it reads the outcome from the course's external approval system through an adapter and shows when the last read succeeded. In every mode it updates the team-project status through M\_Projects.

\subsection{Uses}

\begin{itemize}
    \item M\_HardwareHiding
    \item M\_Authorization
    \item M\_CourseManagement
    \item M\_Projects
    \item M\_StudentsTeams
    \item M\_Notifications
    \item M\_Persistence
    \item M\_AuditLog
\end{itemize}

\subsection{Syntax}

\subsubsection{Exported Constants}
None.

\subsubsection{Exported Access Programs}

\begingroup
\small
\setlength{\tabcolsep}{4pt}
\renewcommand{\arraystretch}{1.2}
\begin{longtable}{>{\raggedright\arraybackslash}p{0.24\textwidth}
                  >{\raggedright\arraybackslash}p{0.34\textwidth}
                  >{\raggedright\arraybackslash}p{0.16\textwidth}
                  >{\raggedright\arraybackslash}p{0.16\textwidth}}
\toprule
\textbf{Name} & \textbf{In} & \textbf{Out} & \textbf{Exceptions} \\
\midrule
\endfirsthead
\toprule
\textbf{Name} & \textbf{In} & \textbf{Out} & \textbf{Exceptions} \\
\midrule
\endhead
\endfoot
\texttt{submit} & user: User, application: ProjectApplication, content: FieldValues & ApprovalRequest & Unauthorized, FeatureDisabled, WrongMode, InvalidStatus, RequestActive, PreflightFailed \\
\texttt{submitRevision} & user: User, application: ProjectApplication, content: FieldValues & ApprovalRequest & Unauthorized, FeatureDisabled, WrongMode, InvalidStatus, RequestActive, PreflightFailed \\
\texttt{recordDecision} & user: User, request: ApprovalRequest, decision: Action, comment: String & -- & Unauthorized, FeatureDisabled, WrongMode, NotActive, InvalidDecision \\
\texttt{recordExternalDecision} & user: User, application: ProjectApplication, decision: Action, comment: String & -- & Unauthorized, FeatureDisabled, WrongMode, InvalidStatus, InvalidDecision \\
\texttt{synchronize} & application: ProjectApplication & -- & WrongMode, InvalidStatus, ExternalServiceError \\
\texttt{getRequest} & user: User, request: ApprovalRequest & ApprovalRequest & Unauthorized \\
\texttt{getPendingRequests} & user: User, course: Course & $\operatorname{List}\langle\text{ApprovalRequest}\rangle$ & Unauthorized \\
\texttt{getApprovedContent} & user: User, application: ProjectApplication & $\text{FieldValues}\cup\{\text{None}\}$ & Unauthorized \\
\texttt{getSyncStatus} & user: User, course: Course & SyncStatus & Unauthorized \\
\bottomrule
\caption{Exported access programs of the Approval Workflow module}
\label{TblApprovalWorkflowAccessPrograms}
\end{longtable}
\endgroup

\subsection{Semantics}

\subsubsection{State Variables}

\begin{itemize}
    \item $\mathit{requests}:\operatorname{Map}\langle\text{String},\text{ApprovalRequest}\rangle$ - in-platform approval requests indexed by identifier, initially loaded from M\_Persistence.

    \item $\mathit{approved}:\operatorname{Map}\langle\text{String},\text{FieldValues}\rangle$ - the approved content currently in force, indexed by application identifier, initially loaded from M\_Persistence.

    \item $\mathit{syncStatus}:\operatorname{Map}\langle\text{Course},\text{SyncStatus}\rangle$ - the latest read of each Integrated course's external system, initially loaded from M\_Persistence.
\end{itemize}

\subsubsection{Environment Variables}

The external approval system of a course (for example GitLab), reached through the adapter named in the course settings over server network connections (through M\_HardwareHiding), the system clock, and persistent request and approval records.

\subsubsection{Assumptions}

\begin{itemize}
    \item Users are authenticated through M\_Authentication before any routine is called.
    \item The caller has re-verified the user, as required by \texttt{FR\_ReverifyUsers}, before calling \texttt{recordDecision} or \texttt{recordExternalDecision}.
    \item The course settings hold \texttt{approvalMode} (ApprovalMode), \texttt{externalAdapter} (used only in mode Integrated), and \texttt{submissionRequirements} ($\operatorname{Map}\langle\text{String},\text{bool}\rangle$, the course's submission items, each mapped to true if required). They are read with $\operatorname{getSettings}$ and changed only through M\_CourseManagement.
    \item In modes External and Integrated the team submits to the external system, so \texttt{submit} and \texttt{submitRevision} are available only in mode InPlatform. Changes to an approved project in the other modes are reviewed outside Keystone and have no request here.
    \item A decision is made by a person (an instructor or TA, as \texttt{hasAccess} decides), never automatically. In mode Integrated the adapter only reads a decision that a person made in the external system.
    \item \texttt{synchronize} is called only by a trusted scheduled job of Keystone.
    \item $\operatorname{notify}(event,record)$ is the M\_Notifications routine that decides the recipients, stores the notification in the caller's transaction, and delivers it later.
\end{itemize}

\subsubsection{Access Routine Semantics}

\paragraph{\texttt{submit(user, application, content)}}
\begin{itemize}
    \item transition: Let $c$ be the course of the application's project. If
    \[
    \begin{aligned}
    &\operatorname{hasAccess}(user,c,\text{Submit},application)=\text{true}
    \land\operatorname{isFeatureEnabled}(user,c,\text{ProjectApprovals})=\text{true}\\
    &\land\operatorname{mode}(user,c)=\text{InPlatform}
    \land application.status\in\{\text{Agreed},\text{TeamProjectChangesRequested}\}\\
    &\land\text{no active request exists for }application
    \land\operatorname{preflight}(user,application,content,\text{Initial})=\emptyset,
    \end{aligned}
    \]
    create a request $r$ with $r.application=application$, $r.content=content$, $r.kind=\text{Initial}$, $r.status=\text{TeamProjectUnderReview}$, and $r.comment=\epsilon$:
    \[
    \mathit{requests}'=\mathit{requests}\cup\{r.id\mapsto r\}.
    \]
    Call $\operatorname{updateApplicationStatus}(application,\text{TeamProjectSubmitted})$ and then $\operatorname{updateApplicationStatus}(application,\text{TeamProjectUnderReview})$ (the checks have passed, so the submission is forwarded for review). Call $\operatorname{recordPendingWork}(c,\text{ProjectApprovals},r.id,\text{true})$ and $\operatorname{notify}(\text{TeamProjectSubmitted},r)$. All changes are saved in one transaction and audited.
    \item output: $out=r$.
    \item exception:
    \begin{itemize}
        \item Unauthorized if $\operatorname{hasAccess}(user,c,\text{Submit},application)=\text{false}$.
        \item FeatureDisabled if ProjectApprovals is not enabled in the course.
        \item WrongMode if the mode of $c$ is not InPlatform.
        \item InvalidStatus if $application.status\notin\{\text{Agreed},\text{TeamProjectChangesRequested}\}$.
        \item RequestActive if the application already has an active request.
        \item PreflightFailed if $\operatorname{preflight}$ finds problems. The exception lists the missing or invalid items and the rules the team breaks, and nothing is sent to the reviewers.
    \end{itemize}
\end{itemize}

\paragraph{\texttt{submitRevision(user, application, content)}}
\begin{itemize}
    \item transition: As \texttt{submit}, except that $application.status$ must be Approved, the new request has $r.kind=\text{Revision}$, the application's status becomes ApprovedChangeUnderReview (a single call of $\operatorname{updateApplicationStatus}$), and $\operatorname{preflight}(user,application,content,\text{Revision})$ is used. $\mathit{approved}[application.id]$ is not changed.
    \item output: $out=r$.
    \item exception: As \texttt{submit}, with InvalidStatus if $application.status\ne\text{Approved}$.
\end{itemize}

\paragraph{\texttt{recordDecision(user, request, decision, comment)}}
\begin{itemize}
    \item transition: Let $c$ be the course of the request's application's project. If $\operatorname{mode}(user,c)=\text{InPlatform}$, $\operatorname{isFeatureEnabled}(user,c,\text{ProjectApprovals})=\text{true}$, $request.status=\text{TeamProjectUnderReview}$, $decision\in\{\text{Approve},\text{Reject},\text{RequestChanges}\}$, and $\operatorname{hasAccess}(user,c,decision,request)=\text{true}$: $\operatorname{finalize}(request,decision,comment)$.
    \item output: None.
    \item exception:
    \begin{itemize}
        \item Unauthorized if $\operatorname{hasAccess}(user,c,decision,request)=\text{false}$.
        \item FeatureDisabled if ProjectApprovals is not enabled in the course.
        \item WrongMode if the mode of $c$ is not InPlatform.
        \item NotActive if $request.status\ne\text{TeamProjectUnderReview}$.
        \item InvalidDecision if $decision$ is not Approve, Reject, or RequestChanges.
    \end{itemize}
\end{itemize}

\paragraph{\texttt{recordExternalDecision(user, application, decision, comment)}}
\begin{itemize}
    \item transition: Let $c$ be the course of the application's project. If $\operatorname{mode}(user,c)=\text{External}$, ProjectApprovals is enabled, $application.status\in\{\text{Agreed},\text{TeamProjectSubmitted},\text{TeamProjectUnderReview},\text{TeamProjectChangesRequested}\}$, $decision\in\{\text{Approve},\text{Reject},\text{RequestChanges}\}$, and $\operatorname{hasAccess}(user,c,decision,application)=\text{true}$: $\operatorname{applyOutcome}(application,decision,comment)$. This records the outcome of the course's external process by hand.
    \item output: None.
    \item exception:
    \begin{itemize}
        \item Unauthorized if $\operatorname{hasAccess}(user,c,decision,application)=\text{false}$.
        \item FeatureDisabled if ProjectApprovals is not enabled in the course.
        \item WrongMode if the mode of $c$ is not External.
        \item InvalidStatus if the application's status is not one of those listed.
        \item InvalidDecision if $decision$ is not Approve, Reject, or RequestChanges.
    \end{itemize}
\end{itemize}

\paragraph{\texttt{synchronize(application)}}
\begin{itemize}
    \item transition: Let $c$ be the course of the application's project. If the mode of $c$ is Integrated and $application.status\in\{\text{Agreed},\text{TeamProjectSubmitted},\text{TeamProjectUnderReview},\text{TeamProjectChangesRequested}\}$, let $(d,m)=\operatorname{adapterFetch}(c,application)$. On success, $\mathit{syncStatus}'[c]=(\mathit{now},\text{false})$, and if $d\ne\text{None}$ then $\operatorname{applyOutcome}(application,d,m)$. If the adapter cannot be reached, $\mathit{syncStatus}'[c]=(\mathit{syncStatus}[c].lastSuccess,\text{true})$, the failure is audited, and $\operatorname{notify}(\text{ExternalApprovalError},c)$ is called. The changes are saved in one transaction.
    \item output: None.
    \item exception:
    \begin{itemize}
        \item WrongMode if the mode of $c$ is not Integrated.
        \item InvalidStatus if the application's status is not one of those listed.
        \item ExternalServiceError if the adapter cannot be reached (the failed state is still saved).
    \end{itemize}
\end{itemize}

\paragraph{\texttt{getRequest(user, request)}}
\begin{itemize}
    \item transition: None.
    \item output: If $\operatorname{hasAccess}(user,c,\text{Read},request)=\text{true}$, where $c$ is the course of the request's application's project:
    \[
    out=\mathit{requests}[request.id].
    \]
    \item exception: Unauthorized if $\operatorname{hasAccess}(user,c,\text{Read},request)=\text{false}$.
\end{itemize}

\paragraph{\texttt{getPendingRequests(user, course)}}
\begin{itemize}
    \item transition: None.
    \item output: If $\operatorname{hasAccess}(user,course,\text{Read},course)=\text{true}$, the requests of the course under review that the user may read, oldest first:
    \[
    out=[r\mid r\in\mathit{requests},\ r.status=\text{TeamProjectUnderReview}\land\operatorname{hasAccess}(user,course,\text{Read},r)].
    \]
    \item exception: Unauthorized if $\operatorname{hasAccess}(user,course,\text{Read},course)=\text{false}$.
\end{itemize}

\paragraph{\texttt{getApprovedContent(user, application)}}
\begin{itemize}
    \item transition: None.
    \item output: If $\operatorname{hasAccess}(user,c,\text{Read},application)=\text{true}$, where $c$ is the course of the application's project:
    \[
    out=
    \begin{cases}
        \mathit{approved}[application.id], & application.id\in\operatorname{dom}(\mathit{approved}),\\
        \text{None}, & \text{otherwise}.
    \end{cases}
    \]
    \item exception: Unauthorized if $\operatorname{hasAccess}(user,c,\text{Read},application)=\text{false}$.
\end{itemize}

\paragraph{\texttt{getSyncStatus(user, course)}}
\begin{itemize}
    \item transition: None.
    \item output: If $\operatorname{hasAccess}(user,course,\text{Read},course)=\text{true}$:
    \[
    out=\mathit{syncStatus}[course],
    \]
    with $\mathit{lastSuccess}=\text{None}$ and $\mathit{failed}=\text{false}$ if no read has been made.
    \item exception: Unauthorized if $\operatorname{hasAccess}(user,course,\text{Read},course)=\text{false}$.
\end{itemize}

\subsubsection{Local Functions}

\begin{description}
    \item[mode(user: User, course: Course): ApprovalMode] Returns the \texttt{approvalMode} course setting, read with $\operatorname{getSettings}(user,course)$.
    \item[preflight(user: User, application: ProjectApplication, content: FieldValues, kind: RequestKind): $\operatorname{Set}\langle\text{String}\rangle$] Returns the names of the problems found before a submission is sent, or the empty set if there are none. A problem is reported for each item marked required in the course's \texttt{submissionRequirements} that is missing or empty in $content$; for each rule in $\operatorname{validateTeam}(user,application.team)$; if the application's project is not Open or Full (it is not available); and, when $kind=\text{Initial}$, if the team has another application with status Approved or ApprovedChangeUnderReview (from $\operatorname{getTeamApplications}$).
    \item[adapterFetch(course: Course, application: ProjectApplication): $(\text{Action}\cup\{\text{None}\})\times\text{String}$] Asks the course's external system, through its adapter, for the decision on the team's project and returns the decision (None if none has been made yet) with its comment. Raises ExternalServiceError if this fails.
    \item[applyOutcome(application: ProjectApplication, decision: Action, comment: String): --] For an application in status Agreed or TeamProjectChangesRequested, first calls $\operatorname{updateApplicationStatus}$ to TeamProjectSubmitted, and for TeamProjectSubmitted then to TeamProjectUnderReview. Then calls it once with Approved (Approve), TeamProjectRejected (Reject, which releases the team's slot), or TeamProjectChangesRequested (RequestChanges). Calls $\operatorname{notify}(\text{TeamProjectDecision},application)$ with the comment and audits the change.
    \item[finalize(request: ApprovalRequest, decision: Action, comment: String): --] Stores $comment$ and ends the request, with $p=request.application$:
    \begin{itemize}
        \item Approve: $\mathit{requests}'[request.id].status=\text{Approved}$, $\mathit{approved}'=\mathit{approved}\oplus\{p.id\mapsto request.content\}$ (this replaces any earlier approval and is the only place an approval is replaced), and $\operatorname{updateApplicationStatus}(p,\text{Approved})$.
        \item Reject: $\mathit{requests}'[request.id].status=\text{TeamProjectRejected}$. If $request.kind=\text{Initial}$, call $\operatorname{updateApplicationStatus}(p,\text{TeamProjectRejected})$, which releases the team's project slot. If $request.kind=\text{Revision}$, call $\operatorname{updateApplicationStatus}(p,\text{Approved})$ and $\mathit{approved}$ is unchanged.
        \item RequestChanges: $\mathit{requests}'[request.id].status=\text{TeamProjectChangesRequested}$. If $request.kind=\text{Initial}$, call $\operatorname{updateApplicationStatus}(p,\text{TeamProjectChangesRequested})$, and the team keeps its slot. If $request.kind=\text{Revision}$, call $\operatorname{updateApplicationStatus}(p,\text{Approved})$ and $\mathit{approved}$ is unchanged.
    \end{itemize}
    In every case, call $\operatorname{recordPendingWork}(c,\text{ProjectApprovals},request.id,\text{false})$ and $\operatorname{notify}(\text{TeamProjectDecision},request)$, and audit the change. Either all of these changes are saved or none are.
\end{description}

\newpage

\section[MIS of Threads Module]{MIS of Threads Module (\moduleid{M_Threads})}
\label{M_Threads}

\subsection{Module}

This module manages discussion threads attached to project proposals, project applications, and team formations. It creates a thread on a record that allows threads, adds replies, deletes threads, and lists the threads of a record, and it notifies the people involved of new threads and replies. Which kinds of records can have threads, and how many threads each record may have, are decided inside this module.

\subsection{Uses}

\begin{itemize}
    \item M\_Authorization
    \item M\_Projects
    \item M\_StudentsTeams
    \item M\_Notifications
    \item M\_Persistence
    \item M\_AuditLog
\end{itemize}

\subsection{Syntax}

\subsubsection{Exported Constants}
None.

\subsubsection{Exported Access Programs}

\begingroup
\small
\setlength{\tabcolsep}{4pt}
\renewcommand{\arraystretch}{1.2}
\begin{longtable}{>{\raggedright\arraybackslash}p{0.24\textwidth}
                  >{\raggedright\arraybackslash}p{0.34\textwidth}
                  >{\raggedright\arraybackslash}p{0.16\textwidth}
                  >{\raggedright\arraybackslash}p{0.16\textwidth}}
\toprule
\textbf{Name} & \textbf{In} & \textbf{Out} & \textbf{Exceptions} \\
\midrule
\endfirsthead
\toprule
\textbf{Name} & \textbf{In} & \textbf{Out} & \textbf{Exceptions} \\
\midrule
\endhead
\endfoot
\texttt{createThread} & user: User, course: Course, kind: ThreadTarget, target: Value, subject: String, content: String & Thread & Unauthorized, TargetNotAllowed, ThreadLimitReached, EmptyContent \\
\texttt{addReply} & user: User, thread: Thread, content: String & -- & Unauthorized, EmptyContent \\
\texttt{deleteThread} & user: User, thread: Thread & -- & Unauthorized \\
\texttt{getThreads} & user: User, course: Course, kind: ThreadTarget, target: Value & $\operatorname{List}\langle\text{Thread}\rangle$ & Unauthorized, TargetNotAllowed \\
\texttt{getThread} & user: User, thread: Thread & Thread & Unauthorized \\
\bottomrule
\caption{Exported access programs of the Threads module}
\label{TblThreadsAccessPrograms}
\end{longtable}
\endgroup

\subsection{Semantics}

\subsubsection{State Variables}

\begin{itemize}
    \item $\mathit{threads}:\operatorname{Map}\langle\text{String},\text{Thread}\rangle$ - threads indexed by identifier, initially loaded from M\_Persistence.
\end{itemize}

\subsubsection{Environment Variables}

Persistent thread and reply records. M\_Persistence records when each thread and reply was written.

\subsubsection{Assumptions}

\begin{itemize}
    \item Users are authenticated through M\_Authentication before any routine is called.
    \item M\_Authorization enforces the SRS visibility rules through \texttt{hasAccess}, so a user can create, read, or reply to a thread only where the SRS allows.
    \item The caller has re-verified the user, as required by \texttt{FR\_ReverifyUsers}, before calling \texttt{deleteThread}.
    \item The target passed in is a Project (for ProjectProposal), a ProjectApplication (for Application), or a Team (for Team). The module reads it again through $\operatorname{getProject}$, $\operatorname{getApplication}$, or $\operatorname{getTeam}$, so it works on the current record and only for a user who may read it.
    \item $\operatorname{notify}(event,record)$ is the M\_Notifications routine that decides the recipients, stores the notification in the caller's transaction, and delivers it later.
\end{itemize}

\subsubsection{Local Constants}

$\mathit{threadLimit}:\operatorname{Map}\langle\text{ThreadTarget},\mathbb{N}\cup\{\text{None}\}\rangle$ is the number of threads each record of a kind may have, where None means no limit and 0 means the kind of record cannot have threads. It is fixed in this module and is currently
\[
\mathit{threadLimit}=\{\text{ProjectProposal}\mapsto\text{None},\ \text{Application}\mapsto 1,\ \text{Team}\mapsto\text{None}\}.
\]
These values, and the set of target kinds, are the secret of this module (see \texttt{LC\_ThreadTargets} in the SRS).

\subsubsection{Access Routine Semantics}

\paragraph{\texttt{createThread(user, course, kind, target, subject, content)}}
\begin{itemize}
    \item transition: If $\operatorname{hasAccess}(user,course,\text{Create},\operatorname{current}(user,kind,target))=\text{true}$, $\mathit{threadLimit}[kind]\ne0$, $\operatorname{targetKind}(target)=kind$, $content\ne\epsilon$, and ($\mathit{threadLimit}[kind]=\text{None}\lor\operatorname{threadCount}(kind,target.id)<\mathit{threadLimit}[kind]$), create a thread $h$ with $h.course=course$, $h.kind=kind$, $h.targetID=target.id$, $h.author=user$, $h.subject=subject$, $h.content=content$, and $h.replies$ empty:
    \[
    \mathit{threads}'=\mathit{threads}\cup\{h.id\mapsto h\}.
    \]
    Call $\operatorname{notify}(\text{ThreadCreated},h)$. Persist and audit the changes.
    \item output: $out=h$.
    \item exception:
    \begin{itemize}
        \item Unauthorized if $\operatorname{hasAccess}(user,course,\text{Create},\operatorname{current}(user,kind,target))=\text{false}$.
        \item TargetNotAllowed if $\mathit{threadLimit}[kind]=0$ or $\operatorname{targetKind}(target)\ne kind$.
        \item ThreadLimitReached if $\operatorname{threadCount}(kind,target.id)\ge\mathit{threadLimit}[kind]$.
        \item EmptyContent if $content=\epsilon$.
    \end{itemize}
\end{itemize}

\paragraph{\texttt{addReply(user, thread, content)}}
\begin{itemize}
    \item transition: If $\operatorname{hasAccess}(user,thread.course,\text{SendMessage},thread)=\text{true}$ and $content\ne\epsilon$, a reply with author $user$ and content $content$ is added at the end of $\mathit{threads}'[thread.id].replies$. Call $\operatorname{notify}(\text{ThreadReply},thread)$. Persist and audit the changes.
    \item output: None.
    \item exception:
    \begin{itemize}
        \item Unauthorized if $\operatorname{hasAccess}(user,thread.course,\text{SendMessage},thread)=\text{false}$.
        \item EmptyContent if $content=\epsilon$.
    \end{itemize}
\end{itemize}

\paragraph{\texttt{deleteThread(user, thread)}}
\begin{itemize}
    \item transition: If $\operatorname{hasAccess}(user,thread.course,\text{Delete},thread)=\text{true}$:
    \[
    \mathit{threads}'=\mathit{threads}\setminus\{thread.id\mapsto thread\}.
    \]
    Persist and audit the change.
    \item output: None.
    \item exception: Unauthorized if $\operatorname{hasAccess}(user,thread.course,\text{Delete},thread)=\text{false}$.
\end{itemize}

\paragraph{\texttt{getThreads(user, course, kind, target)}}
\begin{itemize}
    \item transition: None.
    \item output: If $\operatorname{hasAccess}(user,course,\text{Read},\operatorname{current}(user,kind,target))=\text{true}$ and $\operatorname{targetKind}(target)=kind$, the threads on the target that the user may read, in the order they were created:
    \[
    out=[h\mid h\in\mathit{threads},\ h.kind=kind\land h.targetID=target.id\land\operatorname{hasAccess}(user,course,\text{Read},h)].
    \]
    \item exception:
    \begin{itemize}
        \item Unauthorized if $\operatorname{hasAccess}(user,course,\text{Read},\operatorname{current}(user,kind,target))=\text{false}$.
        \item TargetNotAllowed if $\operatorname{targetKind}(target)\ne kind$.
    \end{itemize}
\end{itemize}

\paragraph{\texttt{getThread(user, thread)}}
\begin{itemize}
    \item transition: None.
    \item output: If $\operatorname{hasAccess}(user,thread.course,\text{Read},thread)=\text{true}$:
    \[
    out=\mathit{threads}[thread.id].
    \]
    \item exception: Unauthorized if $\operatorname{hasAccess}(user,thread.course,\text{Read},thread)=\text{false}$.
\end{itemize}

\subsubsection{Local Functions}

\begin{description}
    \item[current(user: User, kind: ThreadTarget, target: Value): Value] Returns the current target record: $\operatorname{getProject}(user,target)$ for ProjectProposal, $\operatorname{getApplication}(user,target)$ for Application, and $\operatorname{getTeam}(user,target)$ for Team.
    \item[targetKind(target: Value): ThreadTarget] Returns ProjectProposal if $target$ is a Project, Application if it is a ProjectApplication, and Team if it is a Team.
    \item[threadCount(kind: ThreadTarget, targetID: String): $\mathbb{N}$] Returns the number of threads $h$ in $\mathit{threads}$ with $h.kind=kind$ and $h.targetID=targetID$.
\end{description}

\newpage


\bibliographystyle {plainnat}
\bibliography {../../../refs/References}

\newpage

\newpage{}

\section{Appendix --- Generative AI Use Disclosure}

\subsection{AI Tools Used}

Claude Sonnet 5.5.

\subsection{How and Where Used}

Claude was used to draft the MIS of the Projects, Students and Teams, Recommendation, Approval Workflow, and Threads modules, and the shared enums and data structures those modules added to Section~\ref{SecDataTypes} (Program, TeamMemberRole, OfferKind, RuleViolation, RequestKind, ThreadTarget, DateTime, Offer, ApprovalRequest, SyncStatus, and Reply, plus the fields added to the Course, StudentProfile, Team, Project, ProjectApplication, and Thread records). It was also used to fill in the introduction, revision history, and notation rows needed by those modules, and to check the drafts against the SRS and the uses hierarchy of the Module Guide.

\subsection{Verification of AI-Generated Content}

\wss{Describe how the team checked the AI-generated material above against the SRS, the Module Guide, and the implementation before including it. The team is responsible for the accuracy, completeness, consistency, and appropriateness of the submitted document, regardless of whether material was generated by a human or an AI system.}

\newpage{}

\section{Appendix --- Reflection}

\wss{Not required for CAS 741 projects}

The information in this section will be used to evaluate the team members on the
graduate attribute of Problem Analysis and Design.

\input{../../Reflection.text}

\begin{enumerate}
  \item What went well while writing this deliverable? 
  \item What pain points did you experience during this deliverable, and how
    did you resolve them?
  \item Which of your design decisions stemmed from speaking to your client(s)
  or a proxy (e.g. your peers, stakeholders, potential users)? For those that
  were not, why, and where did they come from?
  \item While creating the design doc, what parts of your other documents (e.g.
  requirements, hazard analysis, etc), it any, needed to be changed, and why?
  \item What are the limitations of your solution?  Put another way, given
  unlimited resources, what could you do to make the project better? (LO\_ProbSolutions)
  \item Give a brief overview of other design solutions you considered.  What
  are the benefits and tradeoffs of those other designs compared with the chosen
  design?  From all the potential options, why did you select the documented design?
  (LO\_Explores)
  \item (After you have implemented another team's module, which means this
  isn't filled in until after the original deadline). What did you learn by
  implementing another team's module? Were all the details you needed in the
  documentation, or did you need to make assumptions, or ask the other team
  questions?  If your team also had another team implement one of your modules,
  what was this experience like?  Are there things in your documentation you
  could have changed to make the process go more smoothly for when an
  ``outsider'' completes some of the implementation?

  Your reflection should be about concepts learned in your SE education, such as modular decomposition, cohesion, coupling, information hiding, etc.  You should critique the quality and utility of the specification you received.

  You should create a PR for your implemented module.  A link to the PR should be part of your reflection.

\end{enumerate}

\end{document}